
\section{Upper primes I: the reference law \texorpdfstring{$\mu$}{mu} and its spectral bounds}
\label{sec:upper-reference}

\paragraph{What this section proves and how to read it.}
The upper range is $p\equiv3\pmod4$ in
$[167\cdot2^{40},2^{512})$, partitioned into $310$ consecutive
\emph{bands} (intervals sharing constants; no law lives on a band).
This section constructs $\mu_\Blk$ as a five-type completion of a
$72$-vertex graph law and proves (CP1): reflection positivity at
every cut with uniform marginals. The two \emph{principal cuts}
(called \emph{root cuts} below; the other ten are \emph{nonroot cuts})
determine the completion weights. Displays (PR1)--(PR13)
in \cref{subsec:upper-rows} and formulas (F1)--(F6), (F10) in
\cref{subsec:upper-formulas} are \emph{definitions} of band
input formulas, not lemmas; cite the display number, not a lemma
name.
``Gate'' below means a checked inequality; $L_{\rm band}$ is a
block-size bound $\sum\log p$, not $\Theta$. Inputs: Part~IV
spectral criteria. Outputs: $\mu_\Blk$ and caps for
\cref{sec:upper-comparison}. A first reading should follow the band
definition and (CP1) statement; per-band tables are reference
material (see Appendix~E).
The upper range consists of the blocks of class-three primes in
$[\Pthree,\Pfour)=[167\cdot2^{40},2^{512})$; the primes $p\equiv1\pmod4$ below
$\Pfour$ are atoms of the square-clique prefix (\cref{lem:frontier:iidlocal}).
This section builds, for every block $\Blk$ of an upper band, the reference
$\mu_\Blk$: a nonnegative combination of two-layer field laws and their folded
one-layer companions, with weights chosen so that the five-type completion of
every chamber coordinate and the reserved domination hold. We prove that it is a
probability measure, reflection positive at every cut of the family $\mathfrak F$,
square-affine invariant and with uniform one-slot marginals
(\cref{lem:upper-reference-rp}). This is (CP1) of \cref{thm:comparison} for the
upper atoms (\cref{def:cp1}).

There are $310$ bands (\cref{def:record:upper-bands}): $45$ below $2^{48}$, $64$
in each binary order of $[2^{48},2^{52})$, and nine from $2^{52}$ on. Their inputs
are the following.
\begin{itemize}
\item The two root cuts: the Route~B criterion with graded count
 errors and log-convex nonstar coefficients below $2^{52}$
 (\cref{subsec:record:upper_root}), and the completed root from $2^{52}$.
\item The ten nonroot cuts: class-three tests at $2^{47}$ and $2^{48}$
 (\cref{subsec:upper-nonroot}), with observed merge and constant caps for four
 reserved states (\cref{subsec:upper-observed}).
\item The rows: pinned conditional stars and joint $\rm F/\rm E$ rows with a free
 Young allowance (\cref{subsec:upper-rows,subsec:upper-formulas}).
\end{itemize}
A separate \emph{tail datum} (\cref{def:record:upper-bands}(e)) carries the tail
law (\cref{sec:tail-index}): a fixed set of inputs and weights, evaluated at
$2^{64}$ and valid at both prime classes, used only at the tail primes
$p\ge\Pfour$.

The law $\mu_\Blk$ is the $25$-state five-type completion of \cref{lem:five-type}
applied to the uniform strict-square labelling law of a $72$-vertex two-layer
graph: in each of the two chamber coordinates, folds, matchings, merges and
constants are added with weights in powers of $r=b^{-1/2}$, and one-layer copies
placed on the reserved diagonal restore positivity at the layer flip $\iota$. The
weights must dominate the negative spectra of twelve cut operators, and of the
merge and constant marginals of one of them, and the row errors of six conditional
row laws. These enter as rational \emph{band inputs}, evaluated once per band, at
its start or at a point below it. The two root cuts dominate the construction:
their caps fix the size of all the weights. \Cref{subsec:upper-formulas} lists
the formulas for the inputs with the gates under which they hold;
\cref{lem:record:domain} shows that gates checked at one point hold at every
larger class-three prime and for every block of the band.

The next section builds the comparison law on $\mu_\Blk$ and prices it
(\cref{prop:record:upper}); the section on the tail uses the reference with the
tail datum at every prime $p\ge\Pfour$ of either class (\cref{lem:tail-law}); the
closing reads the upper range only through the prices $S_{\rm u}(a)$,
$A_{\rm u}(a)$.

\emph{Conventions.} $p$ is a prime, $r=p^{-1/2}$ and $q=1/p$; on a block $\Blk$
with least prime $b=b_\Blk$, $r=b^{-1/2}$. Every band records a budget
$L_{\rm band}$ (\cref{def:record:upper-bands}(b)), the block length with which
its data are evaluated and its caps inflated; $L_{\rm u}=50000/27$ is the block
length of the local data of the tail datum and of the completed-root caps from
$2^{52}$ (\cref{def:record:parameters}). The slot array, its chamber cuts
$\vartheta_{c,a}$ and the reserved cut $\iota$ are those of \cref{def:slot-array}.
The five types $\rm L,D,S,R,C$ have the field types ${\rm L,F,E,E,E}$, the formal
degrees $0,1,1,2,3$ and the multiplicities $1,3,2,6,2$
(\cref{def:two-coordinate-five-type}); a \emph{state} is a pair
$\mathfrak t=(\mathfrak t_0,\mathfrak t_1)$ of types, of degree
$\deg\mathfrak t=\deg\mathfrak t_0+\deg\mathfrak t_1$ and multiplicity
$m_{\mathfrak t}=m_{\mathfrak t_0}m_{\mathfrak t_1}$, and its \emph{field pair}
$\mathfrak f(\mathfrak t)$ is the unordered pair of the field types of
$\mathfrak t_0,\mathfrak t_1$. For a cut with a positive envelope we write
$B_{\rm top}$ for its coefficient, $\Op\preceq(1+B_{\rm top}/p)I$: for the
Route~A criterion it is the $B^+$ of \cref{thm:spectral-criterion}(d), for
the Route~B criterion the $B_{\rm top}$ of
\cref{thm:record:direct-star}(e); the letter $B$ is the energy factor of the
bracket. A prime class is $\mathfrak c\in\{1,3\}$, meaning $p\equiv\mathfrak c\pmod4$,
as in \cref{thm:spectral-criterion}. $t_{\rm Y}$ is the Young (mean) allowance of a
row (\cref{lem:five-type}(c),(f)); a band input records it as $t_\varphi$.

\subsection{The \texorpdfstring{$310$}{310} bands and their input data}\label{subsec:upper-bands}

\paragraph{Where the upper range starts.}
On a middle block the comparison already changes records of formal degree one,
and its leading prices are of order $q_\Blk^{2a}b^{-1}$. The upper comparison changes only records
of formal degree at least two, so an upper block $\Blk$ with
least prime $b$ has total comparison price
$s_\Blk+a_\Blk=O(q_\Blk^{2a}b^{-3/2})$ with a band-dependent
constant (\cref{prop:record:upper}(iii)). The upper reference is
the more expensive law to define.
Its weights grow with the root caps, and the root caps grow fast as the primes
decrease: $d_{\rm L}$ is about $12847$ at $\Pthree$, $2557$ at $2^{49}$ and $113$
from $2^{64}$ on (\cref{cert:upper-root-caps,tab:data:upper-bands}). The bands of
$[\Pthree,2^{48})$ carry about $60\%$ of the upper symmetric price
(\cref{prop:record:upper}).

The threshold hypothesis (S1) of the root cuts holds above about $2^{46.34}$
(\cref{subsec:record:upper_root}). All local gates are checked from $2^{47}$ on
(\cref{lem:record:domain}); the rows need no gate on the middle density beyond
$t_{\rm Y}H<1$, since the Young allowance is free (\cref{lem:five-type}(f)). At the
root cuts the completed conditional law of the root has a completion-variance
denominator that is positive only above about $2^{49.96}$ for LL20 and $2^{49.30}$
for LL22. The Route~B criterion expands the root kernel around the
product of the half laws instead; it needs only (S1), so it supplies the root caps
on all bands below $2^{52}$ (\cref{cert:upper-root-caps}). The start
$\Pthree=167\cdot2^{40}$ is a design choice (\cref{def:record:parameters}).

\paragraph{Class three.} Every prime of an upper band is $3$ modulo $4$. At these
primes $-1$ is a nonsquare, the strict-square relation is a tournament, the two
Fourier signs of a spectral test give complex-conjugate matrices
(\cref{thm:spectral-criterion}), and the normalized Paley operator satisfies
$\Phat^*=-\Phat$ and $\Phat^2=-I$ on $L^2_0$ (\cref{subsec:spec:actual-stars}).
The data of the bands are needed, and certified, for class three only: the direct
root caps by Route~B, the nonroot and observed caps, the pinned conditional stars
and joint rows, and the comparisons of \cref{sec:upper-comparison}. The laws of
this section are also used at the tail primes $p\ge\Pfour$ of both classes, but
only with the tail datum (\cref{def:record:upper-bands}(e), \cref{lem:tail-law}),
whose Route~A caps and rows are stated for both classes.

\begin{definition}[Upper bands]\label{def:record:upper-bands}
\leavevmode
\begin{enumerate}[label=(\alph*)]
\item \emph{Bands.} The $310$ bands $[X,Y)$ partition $[\Pthree,\Pfour)$:
 \begin{itemize}
 \item $[167\cdot2^{40},168\cdot2^{40})$;
 \item $[j2^{41},(j+1)2^{41})$ for $84\le j<128$;
 \item $[2^b(1+\frac j{64}),2^b(1+\frac{j+1}{64}))$ for $48\le b\le51$ and
  $0\le j<64$;
 \item $[2^b,2^{b'})$ for consecutive $b<b'$ in $\{52,53,\dots,58,60,64,512\}$.
 \end{itemize}
\item \emph{Blocks.} Below $2^{52}$ a band $[X,Y)$ has two cells $[x_j,x_{j+1})$,
 $x_j=X^{1-j/2}Y^{j/2}$; from $2^{52}$ on it has $J=128\log_2(Y/X)$ cells,
 $x_j=X^{1-j/J}Y^{j/J}$. Each cell carries a recorded integer $n_j\ge1$, and the
 primes $p\equiv3\pmod4$ of the cell are grouped into consecutive blocks of $n_j$
 primes (\cref{lem:record:block-charge}(i)). The integers satisfy
 $n_j\log x_{j+1}\le L_{\rm band}$, where the \emph{band budget} $L_{\rm band}$ is a
 recorded integer between $550$ and $886$. Thus every block has
 $\log q_\Blk\le L_{\rm band}<L_{\rm u}$.
\item \emph{Band inputs.} For each field type $\varphi\in\{\rm L,F,E\}$ the band
 has two inputs, positive rationals with $h\ge1$:
 \begin{itemize}
 \item an \emph{ordinary input} $(d_\varphi,f_\varphi,h_\varphi,M_\varphi,t_\varphi)$;
 \item a \emph{reserved input} $(d^{\rm res}_\varphi,f^{\rm res}_\varphi,
  h^{\rm res}_\varphi,M^{\rm res}_\varphi,t^{\rm res}_\varphi)$.
 \end{itemize}
 For each of the $25$ states $\mathfrak t$ it has a \emph{reserved cap}
 $k_{\mathfrak t}>0$. These are upper bounds, rounded upward in the record, for the
 local quantities of \cref{subsec:upper-formulas}, evaluated for blocks of
 class-three primes at least $X$ and logarithmic length at most $L_{\rm band}$:
 \begin{itemize}
 \item $d_\varphi$ bounds the block caps of the two-layer and of the one-layer
  chamber cut that cut the $\rm L$ coordinate of the graph whose other coordinate
  has field type $\varphi$ (LL20 and LL10, FL21 and FL11, EL21 and EL11).
 \item $(f_\varphi,h_\varphi,M_\varphi,t_\varphi)$ are row coefficients of
  \ref{it:record:upper-F5} valid for the two-layer and for the one-layer rows at
  $\varphi$; in particular $M_\varphi\ge H+t_\varphi$ and $t_\varphi H<1$ for both.
 \item The reserved input uses the one-layer quantities alone: LL10, FL11 or EL11
  and the one-layer row.
 \item $k_{\mathfrak t}$ is the block cap $k_{\mathfrak f(\mathfrak t)}$ of the
  layer cut of the field pair $\mathfrak f(\mathfrak t)$ (LL22, FL22, EL22, FF22,
  FE22, EE22). The exceptions are $k_{\rm LR}=k_{\rm RL}=k_{\rm G}$ and
  $k_{\rm LC}=k_{\rm CL}=k_{\rm C}$, the observed merge and constant caps of EL22
  (\cref{prop:observed-caps}).
 \end{itemize}
 The root entries $d_{\rm L}$ and $k_{\rm LL}$ are the caps of
 \cref{subsec:record:upper_root}. The values are recorded with the band
 (\cref{cert:upper-completion}) and summarized in \cref{tab:data:upper-bands}.
\item \emph{Weights.} Rational weights $w_{\mathfrak t},v_{\mathfrak t}>0$, indexed
 by the $25$ states, with $w_{\rm LL}=1$, such that for each chamber coordinate
 $c\in\{0,1\}$ and each type $\mathfrak t'$ of the other coordinate the five
 weights $a_{\mathfrak s}=w_{\mathfrak t}$ ($\mathfrak t_c=\mathfrak s$,
 $\mathfrak t_{1-c}=\mathfrak t'$, $\mathfrak s\in\{\rm L,D,S,R,C\}$) satisfy
 \begin{equation}\label{eq:upper-slice}
  a_{\rm D}\ge d\,a_{\rm L},\qquad a_{\rm S}\ge M\,a_{\rm D},\qquad
  a_{\rm R}\ge f\,a_{\rm D}+h\,a_{\rm S},\qquad a_{\rm C}a_{\rm S}\ge a_{\rm R}^2
 \end{equation}
 with the ordinary input $(d,f,h,M)$ of the field type of $\mathfrak t'$ (forty
 inequalities); the weights $v$ satisfy the same forty inequalities with the
 reserved inputs; and
 \begin{equation}\label{eq:independent-fold-domination}
  v_{\mathfrak t}\ge k_{\mathfrak t}\,w_{\mathfrak t}\qquad
  \text{for all $25$ states $\mathfrak t$.}
 \end{equation}
\item \emph{Tail datum.} The \emph{tail datum} is a second set of inputs and
 weights, evaluated at $2^{64}$ for blocks of primes of either class. It is not
 the datum of the band $[2^{64},2^{512})$, whose inputs are class-three data.
 \begin{itemize}
 \item Its inputs $(d,f,h)$ and caps are integers, valid for blocks of primes at
  least $2^{64}$ of either class and of length at most $L_{\rm u}$
  (\cref{cert:upper-band-inputs}(d)).
 \item Ordinary inputs: $(d,f,h)=(113,258,120)$ for $\rm L$, $(24,30,11)$ for $\rm F$,
  $(28,8,1)$ for $\rm E$.
 \item Reserved inputs: $(12,18,10)$, $(5,5,3)$, $(5,3,1)$.
 \item $M_\varphi=21/20$ and $t_\varphi=1/50$ for every type and both inputs.
 \item Caps $k_{\mathfrak t}=k_{\mathfrak f(\mathfrak t)}$ with $k_{\rm LL}=45$,
  $k_{\rm LF}=9$, $k_{\rm LE}=7$, $k_{\rm FF}=4$, $k_{\rm FE}=3$, $k_{\rm EE}=1$.
 \item Its $50$ weights satisfy (d).
 \end{itemize}
 It is used only by the tail law (\cref{def:tail-law}).
\end{enumerate}
\end{definition}

\begin{remark}[What is a choice]\label{rem:upper-band-choices}
The bands, the cells and the integers $n_j$ are choices, and any recorded choice
with $n_j\log x_{j+1}\le L_{\rm band}$ is valid. Below $2^{52}$ the integers $n_j$
were chosen with $1/(2a\log x_{j+1})$ as a guide; from $2^{52}$ on the
recorded partitions carry the block
charges $\mathcal U_d$ of \cref{prop:record:upper}, computed at
the anchor exponent $85941/125000000$, which the prices then
reprice (\cref{prop:record:upper}(ii)).
The bands are short below $2^{52}$ because the root caps, which fix the size of all
the weights, fall fast with $X$ ($d_{\rm L}$ is about $12847$ at $\Pthree$, $6569$
at $2^{48}$ and $367$ at $2^{52}$), and each band pays the caps of its own start.

Only the $105$ inequalities of (d) matter. The certificate fixes the weights by
processing the states in increasing type order and taking for each weight the
maximum of its lower bounds, and then checks every inequality exactly. Any
positive rational solution, for instance one chosen jointly by optimization, gives
the same conclusions.

Each band evaluates its inputs at points not above $X$
(\cref{subsec:upper-nonroot}, \cref{cert:upper-band-inputs}). Since positivity is
proved block by block and tensorizes factor by factor, different valid input sets
may be used on different blocks. The data of a band are valid at every
class-three prime $p\ge X$, but a band uses them only on $[X,Y)$. Beyond $2^{512}$
only the tail datum is used.
\end{remark}

\begin{remark}[Band endpoints]\label{rem:band-endpoints}
Band endpoints need not form an equally spaced grid. Each band carries its own
complete prime-count, product-budget, restart, endpoint and jump bounds of
\cref{lem:record:block-charge}. The bands have disjoint prime supports that cover
the range, so the positive band charges add. No boundary term is dropped.
\end{remark}

\subsection{The two-layer reference}\label{subsec:upper-two-layer}

\subsubsection{The two-layer graphs}\label{subsec:push:layer72spectral}

\paragraph{How to read the two-letter names.}
In the upper range every graph, cut and record carries a short
code with one letter for each of the two chamber coordinates. As
\emph{field types}, $\rm L$ is the six-object root coordinate,
$\rm F$ its three-object fold half, and $\rm E$ the three-object
matching coordinate with no forward edges; thus $G_{{\rm FE},2}$ is
the two-layer graph with a fold coordinate and a matching
coordinate. A cut name adds the number of layers and the cut: the
last digit is $0$ or $1$ for the chamber cut in coordinate $0$ or
$1$, and $2$ for the layer cut $\iota$; so LL20 is the chamber cut
in coordinate $0$ of the two-layer root, and EL22 is the layer cut
of the two-layer $\rm EL$ graph. The $25$ \emph{states}
$\mathfrak t$ use instead the five pattern types
$\rm L,D,S,R,C$ of \cref{def:patterns} in each coordinate, so the
state LR has the root pattern in the first coordinate and a merge
pattern in the second. Records are usually named by their field
types (FF, FE, LE), states by their pattern types (DD, DS); where
both occur together the text gives both, as in
\cref{prop:forest-records}.

For the field types let $O_{\rm L}=\Zsix$ and $O_{\rm F}=O_{\rm E}=\{0,1,2\}$ be
the object sets, with the following relations and forward relations. In
$\rm L$, $x\to y$ is related if $y\in\{x+1,x-1,\rot x\}$ (the chamber relation
$\Rel$ without equality), and forward if $y=\rot x$. In $\rm F$ the relation and
the forward relation are their restrictions to $\{0,1,2\}$: the related pairs are
$0\to1$, $1\to0$, $1\to2$, $2\to1$ and $0\to2$, and only $0\to2$ is forward. In
$\rm E$ all ordered pairs of distinct objects are related and none is forward.
For field types $\varphi_0,\varphi_1$ and $j\in\{1,2\}$ let
$G_{\varphi_0\varphi_1,j}$ be the oriented graph on
$O_{\varphi_0}\times O_{\varphi_1}\times\Lambda_j$, $\Lambda_2=\{0,1\}$,
$\Lambda_1=\{0\}$, with an edge $(u_0,u_1,\ell)\to(v_0,v_1,\ell')$ whenever both
coordinate relations $u_0\to v_0$, $u_1\to v_1$ hold and at least one of them is
forward; there is no condition on the layers. The \emph{root} is $G_{\rm LL,2}$.

\begin{lemma}[Two-layer geometry]\label{lem:upper-geometry}
\leavevmode
\begin{enumerate}[label=(\alph*)]
\item Every $G_{\varphi_0\varphi_1,j}$ is an oriented graph. The root has $72$
 vertices and $720$ edges, every vertex has $10$ out-neighbours and $10$
 in-neighbours, and it contains every query
 $(u_0,u_1,\ell)\to(\rot u_0,\rot u_1,1-\ell)$; in particular the twelve queries
 \eqref{eq:upper-twelve} and $\mathcal E_6$ are edges of the root.
\item The chamber reflections of a coordinate of type $\rm L$ and the layer swap
 are automorphisms of the graphs; every permutation of $O_{\rm E}$ is an
 automorphism of a graph with a coordinate of type $\rm E$ (exchangeability).
\item The subgraph of $G_{{\rm L}\varphi,j}$ induced on the plus half
 $\{0,1,2\}$ of its first coordinate is $G_{{\rm F}\varphi,j}$, and likewise in
 the second coordinate. The subgraph of $G_{\varphi\varphi',2}$ induced on the
 layer $\ell=0$ is $G_{\varphi\varphi',1}$.
\item For every $\varphi$ and $j$ the graphs $G_{{\rm F}\varphi,j}$ and
 $G_{{\rm E}\varphi,j}$ induce the same graph on the objects $\{0,1\}$ and on the
 objects $\{1,2\}$ of their first coordinate, and neither has an edge inside one
 object.
\end{enumerate}
\end{lemma}

\begin{proof}
(a) If $u\to v$ is forward in $\rm L$, then $v\to u$ is not related:
$u\notin\{\rot u\pm1,\rot^2u\}$, as $\rot u-u=\pm2$; in $\rm F$ the pair $2\to0$ is
not related. An edge needs a forward coordinate, whose reverse is not related, so no
pair carries edges in both directions; and no loop occurs, since related objects are
distinct.
At a vertex of the root there are $3\cdot3$ pairs of related successors, of which
$2\cdot2$ have no forward coordinate; so there are $5$ out-neighbours in each
layer, $10$ in all, and by the same count $10$ in-neighbours; so $72\cdot10=720$
edges. Both coordinates of a query are forward. (b) A chamber reflection
$\vartheta_a$ commutes with $\rot$ (\cref{lem:slot-array}(i)) and maps $x\pm1$ to
$\vartheta_a(x)\mp1$; the relations do not involve the layer; and $\rm E$ has all
distinct pairs and no forward pair.
(c) Inside $\{0,1,2\}$ the only forward pair of $\rm L$ is $0\to2$, since
$\rot1=5$ and $\rot2=4$; and edges carry no layer condition. (d) Between the
objects $0,1$ and between $1,2$ the relations of $\rm F$ and $\rm E$ are the
same (both directions, none forward); inside one object the relation fails.
The script \path{research/push_2026_10_17/sparse/layer72.py} reconstructs these
facts for all eighteen graphs.
\end{proof}

The spectral caps concern twelve cuts, and the merge and constant marginals of
one of them (\cref{subsec:upper-observed}). The name of a cut lists the field types
of its graph, its number of layers and the cut: $0$ and $1$ for the chamber cut
$\vartheta_{c,0}$ of the coordinate $c$, and $2$ for the layer cut $\iota$. The
two \emph{root cuts} are LL20 and LL22. The ten \emph{nonroot cuts} are the
chamber cuts LL10, FL21, EL21, FL11, EL11 of the $\rm L$ coordinate and the layer
cuts FL22, EL22, FF22, FE22, EE22; coordinate interchange and the chamber
automorphisms give all other cuts. By \cref{lem:upper-geometry}(c) the half graph
of a chamber cut of $G_{{\rm L}\varphi,j}$ is $G_{{\rm F}\varphi,j}$, which is the
graph of the aligned fold, and the half graph of a layer cut is the one-layer
graph. Each cut, with a deletion order fixed once, is spectral data in the sense
of \cref{def:spec:data}; we write $\Op$ for its normalized cut operator on
$L^2(\mu)$, where $\mu=\lambda_p(G^+)$ is the half law, the graph law of the half
graph $G^+$ (\cref{lem:spec:cutkernel}); the actual half marginal is $\varrho\,\mu$,
with $\varrho=\Op\one$. The root cuts have $n=36$ half
vertices, maximal degree $\Delta=20$, internal degrees $d^0_k\le10$ and cross
degrees $r_k\le12$ (LL20) and $r_k\le10$ (LL22).

\subsubsection{The \texorpdfstring{$25$}{25} summand types}\label{subsec:layer72:upper}
For a block $\Blk$ of odd primes and field types $\varphi_0,\varphi_1$ let
$\lambda_{\varphi_0\varphi_1}=\lambda_\Blk(G_{\varphi_0\varphi_1,2})$ and
$\lambda^{(1)}_{\varphi_0\varphi_1}=\lambda_\Blk(G_{\varphi_0\varphi_1,1})$ be the
two-layer and the one-layer field laws (\cref{def:graphs:labeling}): at every
prime of $\Blk$ independently, the uniform law on injective strict-square
labellings. (In this section \emph{pattern summand} below is
what is meant by a summand of $\pi_{\mathfrak t}$ or of
$\mu_\Blk$; \emph{component} keeps its glossary meaning as a
summand of $\nu_\Blk$.) These are the field laws of
\cref{def:two-coordinate-five-type},
with the layer set $\Lambda=\{0,1\}$, respectively $\{0\}$. For a state
$\mathfrak t$ and patterns $\phi=(\phi_0,\phi_1)\in P_{\mathfrak t_0}\times
P_{\mathfrak t_1}$, the \emph{pattern summand} $\Pi_{\mathfrak t,\phi}$ (respectively
$\Pi^{(1)}_{\mathfrak t,\phi}$) is the image of the two-layer (one-layer) field
law of the field types of $\mathfrak t$ under the emission $e_\phi$; objects
outside the image of $\phi$ are forgotten (the \emph{hidden} objects of merges
and constants). The two-layer and one-layer components are written into the
slots $(u_0,u_1,\ell)$ and, for one-layer components, copied onto both layers by
$\operatorname{Diag}_\iota$. Put
\[
 \pi_{\mathfrak t,2}=\sum_\phi\Pi_{\mathfrak t,\phi},\qquad
 \pi_{\mathfrak t,1}=\sum_\phi\Pi^{(1)}_{\mathfrak t,\phi},
\]
of mass $m_{\mathfrak t}$. The \emph{rows} of the coordinate $c$ at the field
type $\varphi$ are the objects of the coordinate $c$ of a field law whose other
coordinate has type $\varphi$; their law is a five-type datum
(\cref{def:five-type-datum}) whose conditional reference is the law of one outer
row given the middle row under the $\rm E$ law (\cref{subsec:upper-rows}). At
$\iota$ the half graph of a two-layer graph is its one-layer graph
(\cref{lem:upper-geometry}(c)); this is why one-layer copies on the reserved
diagonal can cancel the negative envelopes at $\iota$ (\cref{rem:fold-pairing}).

\subsubsection{The reference law \texorpdfstring{$\mu_\Blk$}{mu\_B} and its weights}\label{subsec:independent-fold-weights}

\begin{definition}[The upper reference]\label{def:upper-reference}
Let $\Blk$ be a block of an upper band with least prime $b$ and $r=b^{-1/2}$, and
let $w_{\mathfrak t},v_{\mathfrak t}$ be the weights of the band
(\cref{def:record:upper-bands}(d)). The \emph{reference} of $\Blk$ is
$\mu_\Blk=\widetilde\mu_\Blk/P_\mu(r)$, where
\[
 \widetilde\mu_\Blk=\sum_{\mathfrak t}w_{\mathfrak t}r^{\deg\mathfrak t}\,\pi_{\mathfrak t,2}
 +r\sum_{\mathfrak t}v_{\mathfrak t}r^{\deg\mathfrak t}\operatorname{Diag}_\iota(\pi_{\mathfrak t,1}),
\]
\begin{equation}\label{eq:independent-fold-polynomial}
 P_\mu(r)=\sum_{\mathfrak t}m_{\mathfrak t}\bigl(w_{\mathfrak t}r^{\deg\mathfrak t}
 +v_{\mathfrak t}r^{\deg\mathfrak t+1}\bigr)\ge1 .
\end{equation}
\end{definition}

The polynomial $P_\mu$ has constant coefficient $w_{\rm LL}=1$, the weight of the
root $\pi_{\rm LL,2}$; its coefficient of degree $d$ is the total weight of the
records of overall degree $d$. The one-layer weights $v_{\mathfrak t}$ form an
array of their own: they satisfy the one-layer completion inequalities with the
one-layer inputs and, separately, the domination
\eqref{eq:independent-fold-domination} by the cap of their own state. Weights
$v_{\mathfrak t}$ proportional to $w_{\mathfrak t}$, with one factor for all states,
would have to dominate the largest cap $k_{\rm LL}$ ($9381.137478$ at $\Pthree$,
$2070.828678$ at $2^{49}$) in every state, while the other caps are below $7.01$.

\subsection{The two principal cuts}\label{subsec:record:upper_root}

Two of the twelve cuts are the full degree-twenty root: the first-coordinate cut
LL20, which gives the ordinary input $d_{\rm L}$, and the layer cut LL22, which
gives the reserved cap $k_{\rm LL}$. Their half graphs have $36$ vertices, each of
degree twenty in the root, and their caps are the largest of all: at $\Pthree$
they are about $12847$ and $9381$ ($2557$ and $2071$ at $2^{49}$), against at most
$25$ for the other cuts. Two routes bound them.

\paragraph{The setting.}
Let $(G,\vartheta,U^+,\mathfrak o)$ be the spectral data of LL20 or LL22, with
$n=|U^+|=36$, cross signs $\varsigma$, cross degrees $r_k$, internal degrees
$d^0_k$ and $\Delta=20$. The threshold hypothesis (S1) of \cref{def:spec:data}
is $e_{20}(p,72)<1$ and $p>4n^2=5184$, with the coarse count error $e_t(p,n)$ of
\eqref{eq:graphs:et}. Let $\mu=\lambda_p(G^+)$ be the half law and $\Op$ the
normalized cut operator on $L^2(\mu)$, with the character expansion
\eqref{eq:record:root-kernel} of the root kernel (\cref{lem:spec:expansion}), its
nonstar part $T_{\rm ns}$, its collision part $T_{\rm coll}$ and the star function
$h_\star$ of \eqref{eq:star-function}. At $p\equiv3\pmod4$ this is the setting of
route~B (\cref{subsec:spec:routeB}): $\Hmat^+=\Hmat$, and the two signs of the Schur
test \eqref{eq:record:schur} give complex-conjugate matrices.

The hard root also permits the sign filter of
\cref{lem:signed-collisions}. On its two cuts the exact pin counts are
\begin{equation}\label{eq:upper-signed-collision-counts}
\begin{array}{c|r|r|r|r}
 &\text{all pins}&\text{compatible}&\sum a_{ij}&\sum\widetilde a_{ij}\\ \hline
 \mathrm{LL20}&1296&168&54583296&2752512\\
 \mathrm{LL22}&1296&324&55738368&39813120
\end{array}
\end{equation}
These integers are obtained directly from the $72$-vertex root and
its two reflections (\cref{cert:astra-local}). They sharpen the collision input while preserving
the accepting graph. The certified band caps below retain their
larger collision assignments, so the filter requires no alteration
of the selected comparison laws or closing table.

\paragraph{Route~B at the two root cuts.}
\Cref{thm:record:direct-star} uses no completed conditional law of the root and no
completion-variance denominator, only (S1). Its Route-B frame numbers
(not those of the middle frame)
\ref{it:record:star-H1} are those of \cref{lem:record:frame-inputs}; its inputs
\ref{it:record:star-H3} are the following.

\begin{lemma}[Root inputs]\label{lem:record:direct-star-inputs}
Let the cut be LL20 or LL22, so that $n=36$, $d^0_k\le10$ and $r_k\le12$, and let
$p\equiv3\pmod4$ satisfy (S1). Then hypothesis \ref{it:record:star-H3} of
\cref{thm:record:direct-star} holds with the following choices.
\begin{enumerate}[label=(\roman*)]
\item $M_\star$ is the quantity of \cref{prop:advance:ordered-star} for a fixed
 total order of $U^+$, with the relative count errors $e_i=e^-_{d^0_i}(n)$ of
 \eqref{eq:graphs:graded}.
\item $z_\star=r^2(\sum_iK_i+|B_{\rm even}|)$, with the $K_i$ of
 \cref{prop:advance:ordered-star} and $B_{\rm even}$ of \cref{lem:star-mean}
 for the signs $\chi(\varsigma_{jk})$ of the class $p\equiv3\pmod4$, and
 $Z_\star=2z_\star$.
\item $N_{\rm ns}=\sum_{k=2}^4C_k(Q)r^k$, where $Q\le p$ is a real number at which
 (S1) holds and $C_k(Q)$ bounds the nonstar coefficients of
 \cref{cor:extraction-leaf}(b) for a refinement of the compressed partition
 (\cref{lem:advance:compressed-partition}) with witnesses of sizes $2$, $3$ and $4$.
 \begin{itemize}
 \item These coefficients are evaluated with the graded error pairs and the
  special values $\mathsf S_{\rm m}(s)$, $\mathsf S_\varrho(s)$ of
  \cref{def:extraction:records} and \eqref{eq:leaf-factors}, and summed exactly
  over the complete nonstar histogram of the cut, under the monotone-signature gate
  \eqref{eq:monotone-signature}.
 \item At an interpolation point $Q$ (a prime size, not an exponent anchor) they are this sum.
 \item Between consecutive interpolation points $P_0\le Q<P_1$ the bound is that of
  \cref{lem:logconvex-coefficients}.
 \end{itemize}
\item $D_{\rm col}$ is the minimum, as in \cref{lem:oct24-collision}(e), of
 $D_1+D_{\ge2}$ (with fixed orders and the tail bound of
 \cref{lem:oct24-collision}(c)), the ordered single-coincidence bound
 $D_{\rm ord}$ and the natural bound $D_{\rm nat}$; any one of them may be used
 instead.
\end{enumerate}
\end{lemma}

\begin{proof}
(i), (ii) Given the other labels, the label $x_i$ is uniform on the values
outside the other $n-1$ half labels that satisfy the $d^0_i$ signed constraints
of its internal neighbours, a set with at least
$p2^{-d^0_i}(1-e^-_{d^0_i}(n))$ elements by \cref{lem:extension}(c). So $\mu$ is a
pinned half law (\cref{def:pinned-law}) with count errors $e_i=e^-_{d^0_i}(n)$,
which are below one by (S1), \ref{it:et:compare} and \ref{it:et:poly}. \Cref{lem:star-mean}, with
the signs $\chi(\varsigma_{jk})$ of the class $p\equiv3\pmod4$, gives
$\norm{h_\star-\E_\mu h_\star}_\mu\le M_\star$ and $|\E_\mu h_\star|\le z_\star$, so
$2|\E_\mu h_\star|\le Z_\star$. These bounds concern $\mu$ itself, not the
completed root.
(iii) is \cref{cor:extraction-leaf}(b), with \cref{lem:logconvex-coefficients}
between anchors; the coefficients are nonincreasing in the prime
(\cref{lem:logconvex-coefficients}), so their value at $Q\le p$ bounds them at $p$.
(iv) is \cref{lem:oct24-collision}(e). Neither (iii) nor (iv) uses the class of $p$.
\end{proof}

The hypothesis (S1) holds only above about $2^{46.34}$: $e_{20}(2^b,72)$ is about
$1.125$, $0.796$ and $0.398$ for $b=46,47,49$. The upper range starts at
$\Pthree\approx2^{47.38}$; below about $2^{46.34}$ the Route~B criterion
does not apply to these root cuts.

\begin{lemma}[Monotone endpoint extension]\label{lem:record:monotone-endpoint}
Fix a root cut, the pivot order, the nonstar partition with its records, the
orders of \cref{lem:oct24-collision}, and rationals $s_\star,g_\star$ that satisfy
\eqref{eq:record:schur} for both signs, for the class $p\equiv3\pmod4$, with the
numbers of \cref{lem:record:frame-inputs} evaluated at $r_0=p_0^{-1/2}$, where
$p_0>4n^2$ is a real number (in the applications a power of two) with
$e_\Delta(p_0,2n)<1$. Suppose that $\mathbf R\succeq0$ and that the numbers of
\cref{lem:record:direct-star-inputs}, with $Q=p_0$ in (iii), give $\eps_Z<1$ and
$g_{\rm root}r_0<1$ at $r_0$. Let $p\ge p_0$ be a prime with $p\equiv3\pmod4$. Then
the numbers of \cref{lem:record:frame-inputs,lem:record:direct-star-inputs} at $p$
satisfy \ref{it:record:star-H1}--\ref{it:record:star-H3} with the same
$s_\star,g_\star$, and
\[
 g_{\rm root}(p)\le g_{\rm root}(p_0),\qquad
 B_{\rm top}(p)\le B_{\rm top}(p_0):=\frac{p_0\,\eps_\varrho(p_0)^2}{1-g_{\rm root}(p_0)p_0^{-1/2}} .
\]
Consequently, for every block $\Blk$ of primes $p\ge X\ge p_0$ with
$p\equiv3\pmod4$ and $\log q_\Blk\le L_{\rm u}$, the normalized block root operator
is at least $-d\,b_\Blk^{-1/2}$, with $d=g_{\rm root}(p_0)\exp\bigl(B_{\rm top}(p_0)
L_{\rm u}/(X\log X)\bigr)$.

More generally, let $s_\star,g_\star$ satisfy \eqref{eq:record:schur} at $p_0$ as
above, with $\mathbf R\succeq0$, and let $X\ge p_0$ be real with (S1) at $X$.
Suppose the numbers of \cref{lem:record:direct-star-inputs} at $r_X=X^{-1/2}$, with
$Q=X$ in (iii), give $\eps_Z<1$ and $g_{\rm root}(X)r_X<1$. Let $\Blk$ be a block of
primes $p\ge X$ with $p\equiv3\pmod4$ and $\log q_\Blk\le L'\le L_{\rm u}$. Then the
normalized block root operator of $\Blk$ is at least $-d\,b_\Blk^{-1/2}$ with
$d=g_{\rm root}(X)\exp\bigl(B_{\rm top}(X)L'/(X\log X)\bigr)$, where
$B_{\rm top}(X)=X\eps_\varrho(X)^2/(1-g_{\rm root}(X)r_X)$.
\end{lemma}

\begin{proof}
Write every quantity as a function of $r\in(0,r_0]$. By
\cref{lem:spectral-monotone}(a) the frame quantities $\eps_V$, $\eps_\times$,
$\eps'_H$, $\theta^\mu_j$, $\hat E^2$ and $\eps_{\ne}$ are admissible
(\cref{lem:majorant}); for $\eps'_H$ this includes the entries $\beta_{kl}$ of
\eqref{eq:spec:beta}, maxima of two branches that equal their first branches (the
proof of \cref{lem:spectral-monotone}).
Hence $\eps_{\rm rem}$, $\eps_{\mathbf B}$, $\eps_{\mathbf W}$, $\eps_{\mathbf{EW}}$,
$\alpha_\star$ and $\beta_\star$ are at most their values at $p_0$. The matrices
$\Delta^E$, $\Hmat$ and $\mathbf R$ depend on $p$ only through its class. So
$\mathbf R+\beta_\star I\succ0$ at $p$, because $\mathbf R\succeq0$ and
$\beta_\star\ge\eps_{\mathbf B}\ge\norm\Hmat_{\rm row}/(p-1)>0$; and
\eqref{eq:record:schur} holds at $p$ by \cref{lem:majorant}(d), since
$g_\star/\alpha_\star$ does not decrease and $\beta_\star$ does not increase.

For the inputs: $M_\star/r$ and $z_\star/r^2$ are admissible
(\cref{rem:pivot-order}), hence so are $(M_\star^2+z_\star^2)/r=r(M_\star/r)^2
+r^3(z_\star/r^2)^2$ and $Z_\star$; $N_{\rm ns}/r=\sum_kC_k(p_0)r^{k-1}$ is a
polynomial in $r$ with nonnegative coefficients; and $D_{\rm col}/r$, for the fixed
orders, is at most its value at $p_0$ (\cref{lem:oct24-collision}(e)). Thus
$\eps_Z$, $\eps_{\rm soft}/r$, $g_{\rm root}$ and $\eps_\varrho/r$ are at most their
values at $p_0$: $\eps_Z<1$, so \ref{it:record:star-H3} holds at $p$;
$g_{\rm root}(p)p^{-1/2}\le g_{\rm root}(p_0)r_0<1$; and
$B_{\rm top}(p)=(\eps_\varrho/r)^2/(1-g_{\rm root}r)\le B_{\rm top}(p_0)$. For the
block, the operator is $\bigotimes_{p\in\Blk}\Op$ with
$-g_{\rm root}(p_0)p^{-1/2}I\preceq\Op\preceq(1+B_{\rm top}(p_0)/p)I$ at every $p\in\Blk$
(\cref{thm:record:direct-star}(c),(e)), and
\cref{lem:spec:tensorcomp} gives the bound, since
$\sum_{p\in\Blk}1/p\le L_{\rm u}/(X\log X)$ (\cref{lem:majorant}(c)).

For the general form, the first part of the proof, at $r=r_X\le r_0$, gives
\ref{it:record:star-H1} and \ref{it:record:star-H2} at $X$ with the same
$s_\star,g_\star$, and \ref{it:record:star-H3} at $X$ is the hypothesis $\eps_Z<1$.
The rest of the proof then runs with $X$ in place of $p_0$ for the inputs: for
$p\ge X$ the numbers of \cref{lem:record:direct-star-inputs}, with the coefficients
$C_k(X)$ held fixed, are at most their values at $X$, so $g_{\rm root}(p)\le
g_{\rm root}(X)$ and $B_{\rm top}(p)\le B_{\rm top}(X)$, and
$\sum_{p\in\Blk}1/p\le L'/(X\log X)$. This is the endpoint form of the
construction: the exact Schur witness of an anchor stays valid at every larger
point because every frame, star-Gram and cross error there is at most its anchor
value; this holds by admissibility, and the certificate also checks it separately
for each error (\cref{cert:upper-root-caps}).
\end{proof}

\begin{certificate}[Route~B star-synthesis bounds and Schur tests at the root cuts]\label{cert:record:direct-star}
For LL20 and LL22 and $p_0=2^b$, $b\in\{47,48,49,50,51,52\}$:
\begin{enumerate}[label=(\alph*)]
\item (S1) holds, with $e_{20}(p_0,72)\le0.3978$ at $2^{49}$ and $e_{20}(2^{47},72)$
 about $0.796$; all count errors of
 \cref{lem:record:frame-inputs,lem:record:direct-star-inputs} are below one;
 $\mathbf R\succeq0$ by exact elimination.
\item The frame numbers of \cref{lem:record:frame-inputs}, evaluated in outward
 rational arithmetic and rounded upward, and stored rationals $s_\star,g_\star$
 satisfy \eqref{eq:record:schur}. The records contain $32$ such tests for
 $b\ge49$ (two cuts, four endpoints, both prime classes, both signs; the class-one
 tests are those of the variant described after \cref{lem:record:frame-inputs},
 and class three runs through its real $72\times72$ embedding) and $16$ for
 $b\in\{47,48\}$ (two cuts, two endpoints, both classes, both signs), each
 certified positive semidefinite by exact rational elimination at the value of
 $g_\star$ found for its sign; the left side of \eqref{eq:record:schur} is
 nondecreasing in $g_\star$, so both signs hold at the larger value, which is the
 $g_\star$ used. The class-three matrices are those of the construction; since
 their two signs are complex conjugate, one sign per cut and endpoint would
 suffice. At class three, $g_\star=3193/32$ (LL20) and $18199639/500000$ (LL22) at
 $2^{47}$, and $98128393/10^6$ and $35898609/10^6$ at $2^{48}$.
\end{enumerate}
The records are \path{research/roadmap_2026_10_18/item12/direct_star_spectral.json}
($2^{49}$ to $2^{52}$), with the model and geometry exports
\path{direct_star_models_export.json} and \path{direct_star_geometry_export.json}
of the same directory, produced by \path{direct_star_spectral.py}, and
\path{direct_star_lower_endpoints.json} of the same directory ($2^{47}$, $2^{48}$;
four tests per row, status \texttt{PASS\_EXACT\_DIRECT\_STAR\_ROOT}). The exact
eliminations at $2^{49},\dots,2^{52}$ and the frame errors were repeated by an
independent rational implementation, the four class-three tests at $2^{47}$
by \path{research/advance_2026_10_24/upper/lower_p47/native_review.json}, and the
class-three tests at $2^{48},\dots,2^{51}$ again by
\path{research/continuation_2026_10_24/upper_arithmetic/fractional_roots_review.json}. The
class-three anchor rows ($g_\star$, frame errors, $M_\star$, $z_\star$, collision
bound) for $2^{47},\dots,2^{52}$ are \path{direct_star_collision_v2.json} (same
directory); the reviews \path{research/advance_2026_10_24/upper/review_lower_native47.py}
and \path{research/continuation_2026_10_24/upper_arithmetic/review_roots.py} read
them. The records also derive root caps from these tests with the star constants
of both prime classes; the derived class-one caps are not used, since the
bands below $2^{52}$ have no
class-one primes. Floating-point computations only proposed the rationals
$s_\star,g_\star$; the proof uses the exact tests.
\end{certificate}

\begin{certificate}[Root caps below $2^{52}$]\label{cert:upper-root-caps}
\sloppy
For LL20 and LL22 and every band $[X,Y)$ below $2^{52}$, the inputs of
\cref{lem:record:direct-star-inputs} are taken at the band's start $X$. The
nonstar coefficients come from the complete histograms with the graded errors and
special values, at the $19$ anchors $2^{47},\dots,2^{52}$, $167\cdot2^{40}$ and
$j2^{42}$ for $j\in\{40,42,43,44,46,48,50,52,56,60,80,96\}$, and are interpolated
log-convexly between them (\cref{lem:logconvex-coefficients}). The scalar stars
$M_\star$, $z_\star$, $Z_\star$, the ordered coincidence bound $D_{\rm col}$ and the
rank-one and soft terms are evaluated at $X$; the frame witness $(s_\star,g_\star)$
is that of the band's anchor (\cref{cert:record:direct-star}): $2^{47}$ below
$2^{48}$, and $2^b$ on $[2^b,2^{b+1})$ for $b=48,\dots,51$, every frame, star-Gram
and cross error at $X$ being checked to be at most its anchor value. These give
$\eps_Z<1$, and $g$ below is an upper bound for $g_{\rm root}$, rounded upward to
$10^{-8}$, with $gX^{-1/2}<10^{-3}$. By \cref{lem:record:monotone-endpoint} with
$L'=L_{\rm band}$, the block caps are $d=g\exp(B_{\rm top}L_{\rm band}/(X\log X))$,
rounded upward to $10^{-6}$; they are the root entries $d_{\rm L}$ (LL20) and
$k_{\rm LL}$ (LL22) of the band. Selected rows (the record holds all $301$):
\begin{center}
\footnotesize\setlength{\tabcolsep}{3.5pt}
\begin{tabular}{@{}llllllll@{}}
\toprule
& & \multicolumn{3}{c}{LL20} & \multicolumn{3}{c@{}}{LL22}\\
\cmidrule(lr){3-5}\cmidrule(l){6-8}
$X$ & $L_{\rm band}$ & $g$ & $B_{\rm top}<$ & $d_{\rm L}$
 & $g$ & $B_{\rm top}<$ & $k_{\rm LL}$\\
\midrule
$\Pthree$ & $559$ & $12846.88753069$ & $2.52265\cdot10^{8}$ & $12847.187931$
 & $9381.04100550$ & $1.10946\cdot10^{8}$ & $9381.137478$\\
$2^{48}$ & $566$ & $6568.89967376$ & $8.59331\cdot10^{7}$ & $6568.933791$
 & $5072.00168224$ & $3.71450\cdot10^{7}$ & $5072.013069$\\
$2^{49}$ & $578$ & $2556.55783215$ & $2.30412\cdot10^{7}$ & $2556.559613$
 & $2070.82814422$ & $8.51920\cdot10^{6}$ & $2070.828678$\\
$2^{50}$ & $555$ & $1165.55594825$ & $9.20408\cdot10^{6}$ & $1165.556101$
 & $942.50581374$ & $2.71889\cdot10^{6}$ & $942.505851$\\
$2^{51}$ & $566$ & $612.94721642$ & $4.71877\cdot10^{6}$ & $612.947237$
 & $472.87513870$ & $1.13250\cdot10^{6}$ & $472.875143$\\
$\frac{127}{64}2^{51}$ & $577$ & $370.09363192$ & $2.81352\cdot10^{6}$ & $370.093636$
 & $262.47992391$ & $5.79331\cdot10^{5}$ & $262.479925$\\
\bottomrule
\end{tabular}
\end{center}
At $2^{52}$ and above the Route~A caps $367,\dots,113$ and $261,\dots,45$
of \cref{cert:upper-completed-root} are used. The records are:
\begin{itemize}
\item the coefficient tables
 \path{research/further_2026_10_25/upper/special_coefficients/LL20_coefficients.json}
 and \path{LL22_coefficients.json} ($19$ anchors);
\item their review \path{research/further_2026_10_25/upper/special_coefficients_review.json}
 (\texttt{PASS\_INDEPENDENT\_EVEN\_COUNT\_ROOT\_COEFFICIENTS}, refinement
 \texttt{PASS\_INDEPENDENT\_FORCED\_COINCIDENCE\_ROOT\_COEFFICIENTS}): the
 $30645955$ histogram rows summed exactly, $1596$ upward vertex factors, $114$
 directed coefficients, the $80$ monotone-signature gates \eqref{eq:monotone-signature}
 (every step $s\to s+1$ up to $s+1=\Delta=20$) and
 $3024$ per-signature maxima;
\item the root assembly \path{research/advance_2026_10_24/upper/apply_even_roots.py}
 (the formulas of \cref{thm:record:direct-star} with the interpolation), applied by
 \path{research/further_2026_10_25/upper/apply_special_roots.py} and written into
 the field \texttt{even\_root\_evidence} of
 \path{research/further_2026_10_25/upper/upper_dense_refined.json};
\item its review \path{research/further_2026_10_25/upper/special_upper_review.json}
 (\texttt{PASS\_INDEPENDENT\_EVEN\_COUNT\_ROOT\_UPPER}: $602$ fresh root and block
 caps, $1806$ log-convex coefficients);
\item the scalar root data at the band starts ($M_\star$, $z_\star$, $Z_\star$,
 $D_{\rm col}$, the rank-one and soft terms, the frame errors): for the $45$ bands below
 $2^{48}$ the field \texttt{roots} of the band record, computed with the anchor witness
 at $2^{47}$ by \path{research/advance_2026_10_24/upper/lower_prefix.py} and replayed
 by \path{review_lower_native47.py} of the same directory in the regenerated upper chain
 of \path{regenerate_upper_chain.py}, whose $51$ independent reviews the record
 \path{research/advance_2026_10_24/verification/selected.json} lists; for the bands of
 $[2^{48},2^{52})$ \path{research/continuation_2026_10_24/upper_arithmetic/fractional_roots.json},
 reviewed by \path{fractional_roots_review.json} of that directory
 (\texttt{PASS\_INDEPENDENT\_FRACTIONAL\_ROOT\_ASSEMBLY}: $16$ exact anchor tests,
 $1536$ frame dominations). The nonstar values stored with these scalar data come
 from earlier coefficient tables and are not used: the caps take $N_{\rm ns}$ from the
 coefficient tables of the first item.
\end{itemize}
\end{certificate}

\paragraph{Where the root cost lies.} At $\Pthree$ the nonstar term
$N_{\rm ns}/r\approx12715$ is about $99\%$ of $g(\mathrm{LL20})$, while the frame
coefficient is $g_\star=3193/32$, about $99.8$, and the coincidence term
$D_{\rm col}/r$ about $18.6$: the root is expensive because of its many
two-matchings, not because of its frame.

\paragraph{Route A from $2^{52}$: the completed root.}
From $2^{52}$ on, the root caps come from the Route~A criterion
\cref{thm:spectral-criterion}, whose coupling inputs integrate the internal
neighbourhoods of the root through its completed conditional law
(\cref{subsec:spec:completion-variance}). The criterion accepts any admissible values
of its perpendicular and coupling inputs; the root caps use the following ones
(\cref{rem:spectral-criterion:uses}(iii)).

\begin{lemma}[Completed-root caps]\label{lem:upper-completed-root}
Let the cut be LL20 or LL22, let $p_0$ be a real number with (S1) at $p_0$, and let
$\mathfrak c\in\{1,3\}$. At every prime $p\ge p_0$ with $p\equiv\mathfrak c\pmod4$ the
following are inputs of \cref{thm:spectral-criterion}, with the admissibility that
it requires.
\begin{enumerate}[label=(\roman*)]
\item \emph{Perpendicular inputs.} $N_{\rm ns}=\sum_{k=2}^4C_k(Q)r^k$ with the
 nonstar coefficients of \cref{cor:extraction-leaf}(b) at a fixed $Q\le p_0$ at which
 (S1) holds;
 $D_{\rm col}=\min(D_{\rm ord},D_{\rm nat})$, a minimum as in
 \cref{lem:oct24-collision}(e); $M_\star$ the larger over the two classes of the bound of
 \cref{prop:advance:ordered-star}; and $\eps_Z$ the smaller of
 $2(P_{\rm st}+W_{\rm st})+N_{\rm ns}+D_{\rm nat}$ and $Z_\star+N_{\rm ns}+D_{\rm col}$,
 with the standard star constants $P_{\rm st},W_{\rm st}$ of \cref{lem:spec:stars}
 and $Z_\star$ twice the larger over the two classes of $z_\star$
 (\cref{lem:star-mean}).
\item \emph{Coupling inputs.} $R_+$ of \cref{prop:spec2:mucentre}\ref{it:centre:sup},
 formed with this $\eps_Z$; $\eps_H=\min(\eps'_H+\eps_X,\eps^{\rm loc}_H)$
 (\cref{lem:push:conditionalgram}); $\Gamma_E=\min(\hat E^2,\Gamma^{\rm loc}_E)$
 (\cref{lem:projection-diagonal}), the relative completion variances being bounded
 by \cref{lem:push:localvariance}, whose two-variable bounds use the coefficient
 triples of \cref{thm:dl:H2} (\cref{rem:dl:H2-sharp}(ii)); and, for $\mathfrak c=1$,
 $\bar\zeta$ the minimum of \cref{lem:spec2:parity-transfer}, with
 $\eps_\varrho=R_Z(M_\star+N_{\rm ns}+D_{\rm col})$.
\end{enumerate}
Consequently, if at $r_0=p_0^{-1/2}$ these inputs give $\eps_Z<1$, $\bar\eps_V<1$
and $\ell^{\rm loc}_k>0$ for every $k\in U^+$, and the tests
\eqref{eq:spectral-test} of the class $\mathfrak c$ hold for both signs, then the
conclusions of \cref{thm:spectral-criterion} hold at every prime $p\ge p_0$ with
$p\equiv\mathfrak c\pmod4$; and if $g\ge g_{\mathfrak c}$ and $gr_0<1$, then for every
block $\Blk$ of such primes, all at least $X\ge p_0$, with $\log q_\Blk\le L_{\rm u}$,
the normalized block root operator is at least $-d\,b_\Blk^{-1/2}$, with
$d=g\exp(B_{\rm top}L_{\rm u}/(X\log X))$ and $B_{\rm top}=p_0\eps_\varrho(r_0)^2/(1-gr_0)$.
\end{lemma}

\begin{proof}
(i) $N_{\rm ns}$ bounds $\norm{T_{\rm ns}}$ by \cref{cor:extraction-leaf}(b), and
$D_{\rm col}$, $D_{\rm nat}$ bound $\norm{T_{\rm coll}}$ by
\cref{lem:oct24-collision}(e). By \cref{lem:star-mean} with
\cref{prop:advance:ordered-star}, for the signs of the class of $p$ (the half law is
a pinned half law, as in the proof of \cref{lem:record:direct-star-inputs}),
$\norm{h_\star-\E_\mu h_\star}\le M_\star$ and $2|\E_\mu h_\star|\le Z_\star$; the
maxima over the classes cover both. With the standard star inputs
(\cref{prop:spec2:mucentre}\ref{it:centre:standard}) both forms of $\eps_Z$, and so
their minimum, are admissible in \cref{prop:spec2:mucentre}\ref{it:centre:inputs}.
(ii) These are admissible inputs of the coupling block \cref{cor:spec2:coupling}, by
the lemmas cited. Finally $N_{\rm ns}/r=\sum_kC_k(Q)r^{k-1}$ is a polynomial with
nonnegative coefficients, $D_{\rm col}/r$ and $D_{\rm nat}/r$ are admissible by
\cref{lem:oct24-collision}(e), $M_\star/r$ and $z_\star/r^2$ by \cref{rem:pivot-order},
and $P_{\rm st}/r^2$, $W_{\rm st}/r$ are constants (\cref{lem:spectral-monotone}(c)); so
both forms of $\eps_Z$, and their minimum, are admissible (\cref{lem:majorant}(ii)).
$\Gamma_E$, $\eps_H$ and $\bar\zeta$ are admissible by
\cref{lem:spectral-monotone}(a),(c),(d): the count errors $\epsilon^\pm_t$ of
\cref{subsec:spec:completion-variance} are polynomials in $r$ with nonnegative
coefficients and at most the coarse errors $e_t$, so the denominators of
\cref{lem:push:localvariance}(b) are positive by (S1), and $\ell^{\rm loc}_k>0$ at $r_0$.
So \cref{thm:spectral-criterion} applies, and
\cref{cor:spectral-criterion-blocks}, with $P=X$, gives the block bound.
\end{proof}

\begin{certificate}[Completed-root caps]\label{cert:upper-completed-root}
For LL20 and LL22 at every band start $p_0=X$ from $2^{52}$ on, with the inputs of
\cref{lem:upper-completed-root}, the evaluator asserts $\eps_Z<\frac12$,
$\bar\eps_V<1$, $\ell^{\rm loc}_k>0$ for every $k$ and $gr_0<1$; the matrix tests
\eqref{eq:spectral-test} hold for both prime classes and both signs (four exact
rational tests per cut and threshold); and the recorded numbers give the block caps
$d=g\exp(B_{\rm top}L_{\rm u}/(X\log X))<g+1.2\cdot10^{-5}$:
\begin{center}
\small
\begin{tabular}{@{}lllll@{}}
\toprule
$X$ & $g(\mathrm{LL20})$ & $\lceil d\rceil$ & $g(\mathrm{LL22})$ & $\lceil d\rceil$\\
\midrule
$2^{52}$ & $366.42727119$ & $367$ & $260.84850454$ & $261$\\
$2^{53}$ & $249.23718696$ & $250$ & $157.02296656$ & $158$\\
$2^{54}$ & $186.71731661$ & $187$ & $101.54007395$ & $102$\\
$2^{55}$ & $157.47650343$ & $158$ & $76.79680847$ & $77$\\
$2^{56}$ & $140.27191683$ & $141$ & $63.35956145$ & $64$\\
$2^{57}$ & $130.28930487$ & $131$ & $56.17027922$ & $57$\\
$2^{58}$ & $123.91150306$ & $124$ & $51.86566728$ & $52$\\
$2^{60}$ & $117.04430410$ & $118$ & $47.52728217$ & $48$\\
$2^{64}$ & $112.38975257$ & $113$ & $44.77917248$ & $45$\\
\bottomrule
\end{tabular}
\end{center}
At $2^{64}$ the recorded $g$ is attained by the class-one tests. The
Route~A caps of LL20 and LL22 are $3458$ and $1123$ at $2^{50}$ and $723$ and $543$ at
$2^{51}$; they are larger than the class-three caps by Route~B and are not used. The records are the grids
\path{upper_low_grid.json} ($2^{52}$), \path{upper_low_odd_grid.json} ($2^{53}$),
\path{upper_odd_grid.json} ($2^{55}$, $2^{57}$) and \path{upper_augmented_grid.json}
($2^{54}$, $2^{56}$, $2^{58}$, $2^{60}$, $2^{64}$) in
\path{research/advance_2026_10_17/sparse/}, produced by \path{build_upper_low_grid.py},
\path{build_upper_low_odd_grid.py}, \path{build_upper_odd_grid.py} and
\path{build_augmented_grid.py} through \path{augmented_spectral_core.py} and
\path{conditional_star_rank_spectral.py}; the band map
\path{research/roadmap_2026_10_18/item15/UPPER_PROVIDER_COVERAGE.json} records which
grid supplies each band and cut, and \path{replay_upper_psd.py} of the same
directory recomputes the coefficients and repeats every positivity test by
fraction-free elimination. The root entries of the nine bands from $2^{52}$ on, in
\path{research/further_2026_10_25/upper/upper_dense_refined.json}, are these
integers, and so are those of the tail datum ($113$ and $45$, in
\path{research/roadmap_2026_10_18/item12/upper_v2_local.json}).
\end{certificate}

\subsubsection{Cross-weighted root refinement and the final reference}
\label{subsec:upper-cross-weighted}

The final construction reduces the root bounds in the preceding $301$
bands. The resulting smaller completion weights lower the comparison
prices throughout those bands. The improvement uses
the conditional-density cancellation in
\cref{lem:matching:extractionk}. It keeps their physical labelling
laws, complete character partitions, frames and collision bounds.
Put $\rho_t=(1+e_t^-(36))/(1-e_t^-(36))$.
The LL20 half graph has $36$ vertices and $168$ edges; twelve
vertices have internal degree $8$ and twenty-four have internal
degree $10$. The LL22 half graph has $180$ edges, and all $36$
vertices have internal degree $10$.
Each integrated set has at most $34$ vertices and sum of insertion
degrees at most the number of internal edges.

\begin{lemma}[Uniform extraction divisors]\label{lem:upper-cross-divisors}
The common factor $R_{\rm full}^2$ in the retained LL20 and LL22
nonstar certificates can respectively be divided by
\[
 D_{20}=\rho_8^{23}\rho_{10}^{32}/\rho_0^{17},
 \qquad D_{22}=\rho_{10}^{54}/\rho_0^{16}.
\]
This division is valid also for the interpolated coefficient at each
band start.
\end{lemma}

\begin{proof}
The graded error coefficients are nondecreasing and discretely convex
in $t$. Since $x\mapsto2\operatorname{arctanh}x$ is increasing and
convex, the same holds for $\log\rho_t$. Pad the insertion-degree list
by zeros to length $34$. Moving a unit from a smaller positive degree
to a larger degree below $10$ increases its product. Thus
\[
 \widetilde R_U\le H_{20}=\rho_{10}^{16}\rho_8\rho_0^{17},
 \qquad
 \widetilde R_U\le H_{22}=\rho_{10}^{18}\rho_0^{16},
\]
respectively. Divide the old $R_{\rm full}^2$ by these bounds.
The factors $\rho_0$ are retained, since $e_0^-(36)=36/p$.
For interpolation the common factor $M$ of
\cref{lem:logconvex-coefficients} becomes $M/D$, with the same
signature sum. Dividing both anchor coefficients by their respective
$D$ leaves their coefficient-to-common-factor ratios unchanged.
Hence the new interpolant is exactly the old one divided by $D$
at the evaluation point. The new common factor is a product of
nonnegative powers of $\rho_t$ and $(1-\hat e)^{-1}$, so the original
monotonicity and log-convexity proof still applies.
\end{proof}

For each cut and band start $X$, replace its retained nonstar input
$N$ by $N'=N/D(X)$ in the Route~B negative coefficient formula:
\[
 g'=\frac{g_\star+
       (N'+D_{\rm col}+M_\star^2+z_\star^2+\eps_{\rm soft})X^{1/2}}
       {1-Z_\star-N'-D_{\rm col}}.
\]
Here $Z_\star=2z_\star$ is the scalar-star contribution to $\eps_Z$,
$M_\star^2+z_\star^2$ is the rank-one correction, and
$\eps_{\rm soft}$ is the error from the two softened factors in
\cref{thm:record:direct-star}. Their values, and the frame bound
$g_\star$, are those specified in \cref{cert:upper-root-caps}.
Keep the old positive envelope and its
block inflation. If $g,d$ are the old rounded coefficients, a new
block cap is
$d'=10^{-6}\lceil10^6d\,g'/g\rceil$.
This is valid because the old ratio $d/g$ is an upper bound for
the unchanged inflation factor.

Let $d_{\rm new}$ be the larger of the new LL20 cap and the
reserved LL10 cap; let $k_{\rm new}$ be the new LL22 cap. For the
old completed weights define
\[
 h=10^{-12}\left\lceil10^{12}
       \max\left\{\frac{d_{\rm new}}{w_{\rm LD}},
                   \frac{d_{\rm new}}{w_{\rm DL}},
                   \frac{k_{\rm new}}{v_{\rm LL}}\right\}\right\rceil.
\]
If $h\ge1$ retain the old row. Otherwise keep $w'_{\rm LL}=1$,
put $w'_{\mathfrak t}=h w_{\mathfrak t}$ for every other ordinary
state, and put $v'_{\mathfrak t}=h v_{\mathfrak t}$ for all reserved
states. Use $d_{\rm new},k_{\rm new}$ as the new root inputs.
The two ordinary root inequalities and the LL reserved domination
hold by the choice of $h$; each other inequality is homogeneous in
the scaled variables. Thus all $105$ completion and domination
inequalities remain valid. The nonroot inputs and component
marginals are unchanged.

All $301$ bands below $2^{52}$ have $h<1$; the first has
$h=981224812287/10^{12}$. The nine later bands and the far tail are
retained. Every comparison polynomial scales by $h$, since the
ordinary LL record is already accepting. Its full reference
normalizer changes from $P_\mu$ to $1+h(P_\mu-1)$.
The final prices use this new polynomial in
\cref{lem:four-routes-upper-normalized}, and the upper mean is
recomputed at those same prices. The independent check
\path{research/epsilon_2026_10_29/check_cross_weighted.py}
rebuilds both half graphs, all $602$ root bounds, the rational
scales, all $32\,550$ completion inequalities and all $310$
normalizers.

\subsection{The ten non-principal cuts}\label{subsec:upper-nonroot}

The ten nonroot cuts give the other spectral inputs:
\begin{itemize}
\item $d_{\rm F}$, $d_{\rm E}$ from the two-layer chamber cuts FL21, EL21;
\item the one-layer inputs from LL10, FL11, EL11;
\item the reserved caps from the layer cuts FL22, EL22, FF22, FE22, EE22.
\end{itemize}
Their caps are those of route~A, \cref{thm:spectral-criterion}, with the standard
inputs (\cref{rem:spectral-criterion:uses}(ii)): $N_{\rm ns}$ from
\cref{cor:extraction-compressed} on a fresh compressed partition of the cut,
$D_{\rm col}=D_{\rm nat}$, the natural coincidence bound of
\eqref{eq:collision-natural}, the standard star inputs $M_\star=W_{\rm st}$ and
$Z_\star=2(P_{\rm st}+W_{\rm st})$ (\cref{prop:spec2:mucentre}\ref{it:centre:standard}),
and the coupling block \cref{cor:spec2:coupling} with $R_+$ of
\cref{prop:spec2:mucentre}\ref{it:centre:sup}, $\eps_H=\eps'_H+\eps_X$,
$\Gamma_E=\hat E^2$ and $\bar\zeta=\zeta$. This is formula \ref{it:record:upper-F6}
below.

For the bands the test \eqref{eq:spectral-test} is made for the class
$p\equiv3\pmod4$ only, both Fourier signs, each with its own rational
$\tau^{\rm sp}_{3,\pm}$. The recorded cap is
$g=\max_\pm(\cperp+\tau^{\rm sp}_{3,\pm})/(1-\bar\eps_V)$, rounded upward to
$10^{-8}$ and not to an integer:
\begin{center}
\small
\begin{tabular}{@{}lll@{\qquad}lll@{}}
\toprule
cut & $p_0=2^{47}$ & $p_0=2^{48}$ & cut & $p_0=2^{47}$ & $p_0=2^{48}$\\
\midrule
FL21 & $23.24980857$ & $23.07315227$ & EL21 & $24.95017938$ & $24.6357449$\\
FL22 & $7.00914919$ & $6.95874256$ & EL22 & $5.65086987$ & $5.63739898$\\
FF22 & $3.12377057$ & $3.12364965$ & FE22 & $2.41440785$ & $2.41435434$\\
EE22 & $0.00002884$ & $0.00002784$ & LL10 & $10.17351517$ & $10.05779953$\\
FL11 & $4.10415806$ & $4.10396046$ & EL11 & $4.58309033$ & $4.5829422$\\
\bottomrule
\end{tabular}
\end{center}

\emph{Evaluation points.} The caps are evaluated at $2^{47}$ for the bands below
$2^{48}$, and at $2^{48}$ for every band from $2^{48}$ on, those from $2^{52}$
included. In each band they are inflated at its own start $X$ and budget,
$d=g\exp(B_{\rm top}L_{\rm band}/(X\log X))$, rounded upward to $10^{-6}$
(\cref{cert:upper-band-inputs}(a)): for instance $d=23.249822$ for FL21 at
$\Pthree$ and $23.073161$ at $2^{48}$, and the caps of the field pairs LF and LE
are $7.00915$ and $5.65087$ at $\Pthree$.

\emph{Tail datum.} The tail datum and the tail comparison use the caps of the four
tests (two classes, two signs), at $2^{54}$ and $2^{64}$, with integer block caps.
At $2^{64}$ the class-one tests give the recorded $g$ in every cut except FE22 and
EE22, where the two classes tie (\cref{cert:upper-band-inputs}(d)).

\begin{certificate}[Band inputs]\label{cert:upper-band-inputs}
\sloppy
For every one of the $310$ bands, the numerical inputs required
by \cref{def:record:upper-bands}---the caps of the ten
non-principal cuts, the observed merge and constant caps, and the
row quantities---have been evaluated at the band start and pass
the exact tests named below. The items identify the data file
and independent review for each kind of input.
\begin{enumerate}[label=(\alph*)]
\item \emph{Class-three nonroot caps.} The records are
 \path{research/advance_2026_10_24/upper/lower_p47/class3_caps.json} ($2^{47}$) and
 \path{research/advance_2026_10_24/upper/class3_nonroot_caps.json} ($2^{48}$). Each
 holds the ten cuts with $g$ and $B_{\rm top}$ and $20$ exact PSD tests (two signs
 per cut). They are produced by \path{class3_nonroot.py} and
 \path{extend_lower_caps47.py} of the same directory. The review
 \path{review_class3_nonroot.py} (records \path{class3_nonroot_review.json} and
 \path{lower_p47/class3_review.json}, both
 \texttt{PASS\_INDEPENDENT\_CLASS3\_NONROOT\_UPPER}) rebuilds the matrices from their definitions, checks both signs in the positive-semidefiniteness tests exactly, and recomputes every completion weight of \cref{lem:five-type} for the ten non-principal cuts.
\item \emph{Observed caps} (\cref{prop:observed-caps}). The records are
 \path{observed_caps.json} and \path{lower_p47/observed_caps.json}; the reviews
 \path{observed_caps_review.json} and \path{lower_p47/observed_review.json}, both
 \texttt{PASS\_INDEPENDENT\_OBSERVED\_MARGINAL\_CAPS}, cover the actual
 projections of the eighteen half labels to the twelve and the six observed
 labels, $4$ PSD tests and $2622$ frame pairs. All these files are in
 \path{research/advance_2026_10_24/upper/}.
\item \emph{Rows.} The product-row quantities of \ref{it:record:upper-F3} and
 \ref{it:record:upper-F4} are evaluated at the band's start $X$ with
 $\bar n=L_{\rm band}/\log X$ (field \texttt{field\_primitives} of the band
 record), with the pinned extraction at the bucket $\mathfrak q=0$
 (\path{research/advance_2026_10_24/upper/conditional_fine_parts.json}) and, in
 the correlations, the full star matrices of \cref{prop:conditional-star} with the
 decay weights of its part (c), certified at $j2^{42}$ for even $j$ from $32$ to
 $62$ and at $2^{48}$ (\path{conditional_fine_decay_lower_grid/} and
 \path{conditional_star_fine_p48.json} in \path{research/advance_2026_10_24/local/},
 each reviewed, \texttt{PASS\_INDEPENDENT\_FINE\_PINNED\_CONDITIONAL\_STAR\_GATES}).
 The joint coefficients of \ref{it:record:upper-F5}(a) come from one of $17$
 joint-row endpoint tables, the largest not above $X$:
 \begin{itemize}
 \item eleven in \path{research/further_2026_10_25/upper/joint_endpoints/}: $\Pthree$,
  $j2^{42}$ for $j\in\{56,60,80,96,112\}$, $2^{49}$, $2^{50}$, $2^{51}$, $2^{52}$,
  $2^{60}$, each with its review, indexed by \path{index.json};
 \item six in \path{research/advance_2026_10_24/local/}: the grid tables
  \path{p42_42.json}, \path{p44_42.json}, \path{p46_42.json}, \path{p48_42.json},
  \path{p50_42.json} of \path{conditional_joint_potential_grid/} ($j2^{42}$) and
  \path{conditional_joint_potential_p48.json} ($2^{48}$).
 \end{itemize}
 Every table has the status \texttt{PASS\_CONDITIONAL\_JOINT\_ROW\_PRIME\_GATES} and
 an independent review \texttt{PASS\_INDEPENDENT\_CONDITIONAL\_JOINT\_PIN\_POTENTIAL}.
 The reviews of the band records in \path{research/advance_2026_10_24/upper/} are
 \path{joint_rows_review.json}, \path{joint_pins_upper_review.json} and
 \path{joint_potential_upper_review.json}
 (\texttt{PASS\_INDEPENDENT\_JOINT\_CONDITIONAL\_UPPER\_ADAPTER}),
 \path{conditional_fine_parts_review.json} (the bucket $\mathfrak q=0$),
 \path{fine_conditional_fields_review.json} and \path{signed_row_upper_review.json}.
 The allowances are reviewed by
 \path{research/further_2026_10_25/upper/allowances_review.json} and
 \path{endpoint_allowances_review.json}
 (\texttt{PASS\_INDEPENDENT\_UPPER\_ALLOWANCE\_REOPTIMIZATION}), which recompute $f$
 and $M$ from the stored $\ell_{\rm mid}$, $c_{\rm mar}$, $(k_{\rm J},c_{\rm F})$, $h$
 and $t_\varphi$ by the formula of \ref{it:record:upper-F5}(a) and authenticate every
 pair $(k_{\rm J},c_{\rm F})$ against its table. At $\Pthree$, for instance, the
 ordinary inputs have $(M_\varphi,t_\varphi)=(1.012767561,\,0.012758561)$,
 $(1.023485,\,0.023484)$ and $(1.000058323,\,0.000057323)$ for $\rm L$, $\rm F$, $\rm E$.
\item \emph{Tail datum and tail caps (both classes).} The nonroot caps and the six
 Schur rows of the tail datum are those of the band $[2^{64},2^{512})$ of
 \path{research/frontier_2026_10_17/root/complete_c375_verification.json}, which
 \path{research/roadmap_2026_10_18/item12/upper_v2_local.json} keeps. The nonroot
 caps at $2^{54}$ and $2^{64}$ come from the grid \path{upper_augmented_grid.json} of
 \path{research/advance_2026_10_17/sparse/}, built by \path{build_augmented_grid.py}
 with four exact rational tests per cut and threshold (two classes, two signs),
 each of the form \eqref{eq:spectral-test}, and the gates of
 \ref{it:record:upper-F6}. The replays \path{replay_upper_psd.py} and
 \path{replay_upper_local_inputs.py} of \path{research/roadmap_2026_10_18/item15/}
 recompute every cap, the six rows and their Schur coefficients (with $H<103/100$
 asserted for every row), and the integer inputs.
\end{enumerate}
The searches only chose finite partitions and rational parameters; the proof uses
the exact tests and the stated formulas.
\end{certificate}

\subsection{Merge and constant bounds for four states}\label{subsec:upper-observed}

A merge component of the upper reference depends only on the
labels in rows $0$ and $1$ of its $\rm E$ coordinate, and a
constant component only on row $0$
(\cref{subsec:layer72:upper}). For the order at the reserved cut,
only the compression of the layer-cut operator to functions of
those labels (the \emph{observed} labels) therefore matters, and
\cref{lem:five-type}(e) allows such a component a cap of its own.
For the four states in question---denoted LR, RL (the merges) and
LC, CL (the constants) in the band data, for the field pair
$\rm LE$---this compressed operator has a much smaller negative
part than the full cut operator (whose cap is $5.65$ at $2^{47}$).
The proposition gives the compressed caps.

\begin{proposition}[Observed merge and constant caps]\label{prop:observed-caps}
Let $p_0\in\{2^{47},2^{48}\}$, let $p\ge p_0$ be a prime with $p\equiv3\pmod4$, let
$\Op$ be the normalized cut operator of EL22 at $p$, and let $\Pi_O$ be the
conditional expectation onto the functions of the labels of a set $O$ of half
vertices. For the merge ($O$ the twelve labels of the rows $0,1$ of the $\rm E$
coordinate) and the constant ($O$ the six labels of the row $0$), the compression
$\Op^O=\Pi_O\Op\Pi_O$ on the functions of $x_O$ satisfies
\[
 -g_{\rm G}rI\preceq\Op^O\preceq(1+B_{\rm top}/p)I,\qquad\text{respectively}\qquad
 {-g_{\rm C}rI}\preceq\Op^O\preceq(1+B_{\rm top}/p)I,
\]
where:
\begin{itemize}
\item $(g_{\rm G},g_{\rm C})=(2.268819655,\,0.079761731)$ and
 $B_{\rm top}<39236.89926$ for $p_0=2^{47}$;
\item $(g_{\rm G},g_{\rm C})=(2.245405779,\,0.056393528)$ and
 $B_{\rm top}<39227.259341$ for $p_0=2^{48}$.
\end{itemize}
Here $B_{\rm top}$ is the positive constant of the full cut EL22. Hence for any $L'$ and a block
of such primes, all at least $X\ge p_0$, with $\log q_\Blk\le L'$, the compressed
block operators have negative coefficients at most
$k_{\rm G}=g_{\rm G}e^{B_{\rm top}L'/(X\log X)}$ and
$k_{\rm C}=g_{\rm C}e^{B_{\rm top}L'/(X\log X)}$. These may serve as the caps of the
reserved states LR, RL and LC, CL in \eqref{eq:independent-fold-domination}
(\cref{lem:five-type}(e)).
\end{proposition}

\begin{proof}
\emph{Geometry.} The plus half of the layer cut of $G_{\rm EL,2}$ is the layer
$\ell=0$, with the eighteen labels $(a,b)$, $a\in\{0,1,2\}$, $b\in\Zsix$. The
edges of the graph are $(a,b,\ell)\to(a',\rot b,\ell')$ for $a\ne a'$ and all
$\ell,\ell'$ (\cref{subsec:push:layer72spectral}: the $\rm E$ coordinate has no
forward pair), so the signed cross-incidence matrix is
\[
 D_{(a,b),(a',b')}=\one_{a\ne a'}\bigl(\one_{b'=\rot b}-\one_{b=\rot b'}\bigr),
\]
an antisymmetric matrix, and the internal and cross degrees are four. Every
permutation of the rows $a$ is an automorphism (\cref{lem:upper-geometry}(b)). The
merge and the constant components are images of the field law under the emission
(\cref{subsec:layer72:upper}); at $\iota$ their cut forms are those of the field
law on the functions of $x_O$, that is, of $\Op^O$. Their observed half
marginals are marginals of the
original half law, not uniform laws of smaller graphs.

\emph{Complete star operator.} There are $162$ distinct supports of size at least
two in the cross neighbourhoods. With the actual features of
\cref{subsec:spec:actual-stars} for the half law and the synthesis
$\mathcal F=(V_j,(W_S/r)_S)$, \cref{cor:star-matrix} writes the doubly centred
complete star part of the cut kernel as $r\mathcal FM(\Phat)\mathcal F^*$, with
\[
 M(\Phat)=\begin{pmatrix}D\otimes\Phat&B_\star\\B_\star^{\mathsf T}&0\end{pmatrix},
 \qquad (B_\star)_{j,S}=(-1)^{|S|}\prod_{i\in S}D_{ji};
\]
since $D$ is antisymmetric, the right-star coefficients $\prod_{i\in S}D_{ij}$ are
the entries of $B_\star^{\mathsf T}$. The higher-star constant couplings vanish only
under the double centring; they stay in the constant-coupling bound below.

\emph{Schur tests.} By \cref{rem:realified-gates}, $\norm{M(\pm i)}\le k$ follows
from the rational tests $k^2I-B_\star B_\star^{\mathsf T}\pm ikD\succeq0$. They hold
with $k=6$ for the full matrix, and with $k=2.188902$ for the principal matrix of
the features of the merge's observed labels (twelve singletons and sixty supports;
on the singleton rows $B_\star B_\star^{\mathsf T}=I$). The principal matrix of the
features of the constant's observed labels is zero. All these are finite rational
matrix inequalities, checked by exact $LDL^{\mathsf T}$ elimination.

\emph{Visible and hidden features.} Call a feature \emph{visible} if it is a
function of $x_O$ (a singleton of $O$, or a support contained in $O$, also when its
principal coefficients vanish) and \emph{hidden} otherwise: the merge has $72$
visible and $108$ hidden features, the constant $12$ and $168$. Let
$\mathcal F_{\rm v}$, $\mathcal F_{\rm h}$ be the two parts of the synthesis, so that
$\Pi_O\mathcal F_{\rm v}=\mathcal F_{\rm v}$. By \cref{lem:general-star}(f),
$\norm{\mathcal F_{\rm v}}^2\le\Gamma$ for the actual frame bound $\Gamma$. By
\cref{lem:hidden-feature}, with the best unobserved pivot of each hidden feature and
$D_j\le2^4/(1-e_4(p,18))$ (\cref{lem:extension}), each hidden feature has a bound
for $\norm{\Pi_O\mathcal F_\alpha}$; let $e_{\rm hid}^2$ be the sum of their squares,
so that $\norm{\Pi_O\mathcal F_{\rm h}}\le e_{\rm hid}$ by Cauchy--Schwarz, no
orthogonality of the hidden features being used. Splitting $M$ into its visible
principal block and the three other blocks, each of norm at most $6$, the
compression of the centred star operator has norm at most
\[
 r\bigl\{\Gamma k+6\bigl(2\sqrt\Gamma\,e_{\rm hid}+e_{\rm hid}^2\bigr)\bigr\}.
\]

\emph{Lower bound.} Let $N_{\rm ns}$, $D_{\rm col}$, $R_Z$ and $\eps_\varrho$ be the
nonstar, coincidence, normalizer and constant-coupling bounds of the full cut EL22
at $p$ (\ref{it:record:upper-F6}). The compression $\Pi_O$ contracts the nonstar and
coincidence parts and the constant-coupling vector, and keeps the scalar corner
$\ip\one{\Op\one}=1$. So $A=\Op^O$ has $\norm{P_0A\one}\le\eps_\varrho$ and
$P_0AP_0\succeq-cP_0$ with
\[
 c=R_Z\Bigl[r\bigl\{\Gamma k+6\bigl(2\sqrt\Gamma\,e_{\rm hid}+e_{\rm hid}^2\bigr)\bigr\}
 +N_{\rm ns}+D_{\rm col}\Bigr],
\]
and \cref{prop:matching:positive}(a$'$) gives
$A\succeq-(c+\eps_\varrho^2/(1+c))I\succeq-(c+\eps_\varrho^2)I$. Outward evaluation
of $(c+\eps_\varrho^2)/r$ at $p_0$, rounded upward, gives the stated $g_{\rm G}$ and
$g_{\rm C}$. The upper bound is the compression of the positive envelope
$\Op\preceq(1+B_{\rm top}/p)I$ of the full cut (\cref{thm:spectral-criterion}(d), with
the class-three caps of \cref{cert:upper-band-inputs}(a)).

\emph{All primes and blocks.} Every error bound above, after division by its
indicated power of $r$, is admissible (\cref{lem:majorant}): the extension
denominators increase with $p$, and $\Gamma$, $e_{\rm hid}$, $N_{\rm ns}/r$,
$D_{\rm col}/r$ and $\eps_\varrho^2/r$ decrease, while the rational matrix caps $6$
and $2.188902$ do not depend on $p$; so the bounds hold at every $p\ge p_0$ with the
values at $p_0$. The block statement is \cref{lem:spec:tensorcomp}(b), applied to the
compression to the functions of the labels of $O$, with
 $\sum_{p\in\Blk}1/p\le L'/(X\log X)$ (\cref{lem:spec:tensorcomp}(c)).
\end{proof}

The four caps enter only the reserved dominations
\eqref{eq:independent-fold-domination}, and through them the reserved weights and
the deletion polynomial of the comparison (\cref{def:upper-comparison}). They depend
on the state, not on the placement of the merge or the constant in the slots, and
the weights of \cref{def:record:upper-bands}(d) are symmetric in the two type
coordinates, so the reference $\mu_\Blk$ remains unchanged
when its root (principal) law is replaced by the swapped one
(interchanging the two type coordinates). The block caps
are rounded upward to $10^{-6}$ with the band's own budget: $k_{\rm G}=2.26882$ and
$k_{\rm C}=0.079762$ at $\Pthree$, and $2.245406$ and $0.056394$ from $2^{48}$ on
(\cref{cert:upper-band-inputs}(b)).

\subsection{Conditional laws given the middle label}\label{subsec:upper-rows}

The five-type completion is applied one chamber coordinate at a
time. Fix such a coordinate; its three objects are called the
rows $0,1,2$. The completion needs, after conditioning on all
labels of the middle row $1$, bounds for the conditional laws of
the two outer rows under the fold pattern ($\rm F$) and under the
matching pattern ($\rm E$), measured against one common product
reference (\cref{lem:five-type}(b),(c),(f)); in the upper bands
the two rows are moreover bounded jointly, by one Gram matrix
(\cref{subsec:upper-joint-rows}). In detail: fix a chamber
coordinate $c$, a field type $\varphi$ of the other coordinate and
a layer count $j$. The rows are the objects $0,1,2$ of the
coordinate $c$ in the field laws
$\lambda_\Blk(G_{{\rm F}\varphi,j})$ and
$\lambda_\Blk(G_{{\rm E}\varphi,j})$. Condition on the $z$ distinct
labels of the middle row $1$. By \cref{lem:upper-geometry}(d), in
both laws the two outer rows have $n$ labels, no internal edges,
and exactly the same signed conditions (pins) towards the middle
row; they differ only in the edges between the outer rows
themselves. The six geometries are
\[
\begin{array}{c|cccccc}
 (\varphi,j)&({\rm L},1)&({\rm L},2)&({\rm F},1)&({\rm F},2)&({\rm E},1)&({\rm E},2)\\
 \hline
 n=z&6&12&3&6&3&6\\
 \varpi=\text{maximal pin degree}&2&4&1&2&0&0
\end{array}
\]
These are the pinned rows of \cref{subsec:push:pinnedremainder}, with pin bound
$\varpi$ and $z$ fixed labels, grouped with the bucket $\mathfrak q=0$, so that the
integer free-degree costs $\mathrm{oc}(\mathcal C)$ enter exactly
(\cref{rem:pinned-instances}). Throughout this subsection $r=p^{-1/2}$ for one odd
prime $p$, and the symbols introduced here are local to it; block quantities are
formed in \ref{it:record:upper-F5}.

\subsubsection{The product reference and the star bounds}
\label{subsec:push:layer72product}
Let $e_t=e^-_t(z)$, the graded error of \eqref{eq:graphs:graded}, be the relative
error of a strict $t$-pin count with $z$ excluded labels; assume $e_\varpi<\frac12$,
and put
\[
 \kappa=\frac{2^\varpi}{1-e_\varpi},\qquad I_0=1-\binom n2\kappa r^2 .
\]
An outer label must avoid the middle labels and satisfy its signed pins; its
\emph{allowed set} has at least $p2^{-\varpi}(1-e_\varpi)$ elements
(\cref{lem:extension}(c)), so the uniform law on it has density at most $\kappa$
relative to the uniform law of $\Fp$. Replacing a smaller actual pin degree by $\varpi$
enlarges the density, insertion, denominator and collision bounds of the pinned
extraction below and imposes no edge; the certificates keep the actual pin
dictionaries and check their equality, both orientations, and the absence of
outer-row edges. Let $\lambda$ be the product of the allowed-set laws of the $n$
labels of a row and $\mu$ the law $\lambda$ conditioned on the row being injective
(the \emph{actual row law}); by \cref{lem:restriction}(d) the injection probability
is at least $I_0$. For $k\in\{\rm F,E\}$ let $K_k$ be the raw cross kernel between
the outer rows, relative to $\mu\otimes\mu$: the indicator that the two rows are
disjoint times the product over the cross pairs $(i,i')$ of
$1+\chi(\varsigma_{ii'}(x'_{i'}-x_i))$. The conditional law of the outer rows under
the $k$ law is $K_k\,\mu\otimes\mu$ normalized; put $Z_k=\ip\one{K_k\one}$ and
$\varrho_k=K_k\one/Z_k$, the normalized marginal density.

\emph{One-character means.} For a fixed output $x$ outside the middle labels and an
outer label with allowed law $\mathcal A$, the mean $\E_{\mathcal A}\chi(x-y)$ has
modulus at most
\[
 \beta_\chi:=\frac{\sum_{s=1}^\varpi\binom{\varpi}{s}(s-\one_{s\ \rm odd})r
 +2^\varpi(z+1)r^2}{1-e_\varpi} .
\]
Indeed, expand $2^\varpi\one_{\mathcal A}$ into pin monomials: a monomial with $s\ge1$
pins gives a character sum with the $s+1$ distinct roots $x$ and its pins, of
modulus at most $s\sqrt p$ for even $s$ and $(s-1)\sqrt p+1$ for odd $s$
(\cref{thm:weil}); the monomial without pins gives zero, the constants $1$ of the odd
$s$ add up to $2^{\varpi-1}$, and deleting the $z$ excluded labels costs at most $z2^\varpi$;
divide by $2^\varpi|\mathcal A|\ge p(1-e_\varpi)$. (For $\varpi=4$, $z=12$ this
is $\beta_\chi=(24r+208r^2)/(1-e_4)$.)

\emph{Stars.} Expand each character $\chi(\varsigma(x'-y))$ of the kernel into its
mean and its centred part. Restricting the Paley operator to two allowed sets has
norm at most $\kappa r$, centring an input is contractive, products on distinct input
coordinates have Gram kernels given by Schur products whose extra centred factors
have diagonal at most one, and exact input supports are orthogonal. Consequently,
after its singleton edges are removed, a star of degree $d$ has input-constant part at
most $\mathsf E_d$ and nonconstant part of norm at most $\kappa r\sqrt{\mathsf T_d}$,
where
\[
 \mathsf E_d=(1+\beta_\chi)^d-1-d\beta_\chi,\qquad
 \mathsf T_d=\sum_{s=1}^d\binom ds\bigl((1+\beta_\chi)^{d-s}-\one_{s=1}\bigr)^2
\]
(for a singleton support, removing the singleton edge subtracts one from its
coefficient). Let $(d_i)$ and $(d'_{i'})$ be the two lists of cross degrees and
$n^\times_k$ the number of cross edges of $K_k$. The full star part of $K_k$ has norm
at most
\[
 \mathsf C_k=n^\times_k\kappa r+\sum_i\bigl(\kappa r\sqrt{\mathsf T_{d_i}}+\mathsf E_{d_i}\bigr)
 +\sum_{i'}\bigl(\kappa r\sqrt{\mathsf T_{d'_{i'}}}+\mathsf E_{d'_{i'}}\bigr),
\]
and its centred constant image is at most the larger, over the two orientations, of
\[
 \mathsf M^{\rm st}_k=\kappa r\Bigl(\sum_id_i^2\Bigr)^{1/2}+\sum_i\mathsf E_{d_i}
 +\sum_{i'}\kappa r\sqrt{\mathsf T_{d'_{i'}}} ;
\]
for the singleton part, the centred output coordinates are orthogonal. The single
edges have means at most $\mathsf J_\varpi r^2$, by (PR2) below.

\emph{Injection.} Let $N^{\rm ns}_k$ and $D^{\rm col}_k$ be the nonstar and
cross-coincidence norms of $K_k$ on $L^2(\mu)$: $N^{\rm ns}_k$ is the bound
$N^{\rm pin}$ of \cref{prop:pinned-extraction} for the compressed partition, and
$D^{\rm col}_k$ the bound of \cref{prop:pinned-collision}, whose two orientations
agree because in all twelve row geometries the row and column degree multisets
coincide (\cref{rem:pinned-instances}). By \cref{lem:restriction}(c), applied with
the full star bound $\mathsf C_k$ (a strict-centred test extends to a
product-centred test),
\begin{gather*}
 |Z_k-1|\le\eps_k:=\frac{n^\times_k\mathsf J_\varpi r^2+\sum\mathsf E
 +\mathsf C_k(2\sqrt{1-I_0}+1-I_0)}{I_0}+N^{\rm ns}_k+D^{\rm col}_k,\\
 \norm{P_0K_k\one}\le\mathsf M_k:=\frac{\mathsf M^{\rm st}_k+\mathsf C_k\sqrt{1-I_0}}{I_0}
 +N^{\rm ns}_k+D^{\rm col}_k,
\end{gather*}
where $\sum\mathsf E$ is the sum of $\mathsf E_d$ over both degree lists.
\Cref{cert:upper-band-inputs}(c) asserts $\eps_k<\frac12$.

\subsubsection{Differences of conditional laws given the middle label}\label{subsec:advance:product-rows}
The completion needs the difference of the two conditional laws, not each law
separately: the middle density is a ratio of the normalizers $Z_{\rm F}$, $Z_{\rm E}$,
and the fold marginal densities are ratios $\varrho_{\rm F}/\varrho_{\rm E}$. The
common product reference makes this difference small. The estimates below hold at
every odd prime; the certificate evaluates them for both prime classes and takes the
larger value.

\paragraph{Complete conditional means.}
For distinct-root sets $S,T$ of size at most four,
\[
 \Bigl|\sum_{x,y}\chi\Bigl((x-y)\prod_{s\in S}(x-s)\prod_{s'\in T}(y-s')\Bigr)\Bigr|
 \le40p .\tag{PR1}
\]
Empty root sets give zero, and for $|S|=|T|=1$ the two-root identity
\eqref{eq:two-root} gives modulus at most $p$. Otherwise this is \cref{thm:dl:H2}
(in its symmetric form if $|S|=1$), since $\mathcal B(m,m',s)\le\mathcal B(4,4,0)=
28+8r+4r^2<40$ for $2\le m\le4$, $1\le m'\le4$ (\cref{rem:dl:H2-sharp}(iii)). Expanding the
two allowed-set densities and deleting their $z$ common excluded labels gives, for
independent outer labels $X,Y$ with allowed laws (including $\varpi=0$),
\[
 |\E\chi(X-Y)|\le\mathsf J_\varpi r^2,\qquad
 \mathsf J_\varpi=\frac{40(2^\varpi-1)^2+2z4^\varpi}{(1-e_\varpi)^2} .\tag{PR2}
\]
In the $({\rm L},2)$ geometry the signed $\rm F-E$ singleton matrix $\mathbf D$ has
absolute entry sum $\norm{\mathbf D}_1=72$ in both prime classes, so its full singleton
mean is at most $72\mathsf J_4r^2=1090368r^2/(1-e_4)^2$. All higher-star means, nonstar
remainders and cross-injection corrections are retained. Thus \emph{the conditional
normalizer difference has order $r^2$, not $r$}.

\paragraph{Exact output supports.}
Write $M_{\rm F},M_{\rm E}$ for the signed cross matrices, with zero for an absent edge.
For a centre $j$ and a set of leaves $S$ put
\[
 c_j(S)=\prod_{i\in S}M_{\rm F}(i,j)-\prod_{i\in S}M_{\rm E}(i,j),\qquad
 b_j(T)=\sum_{\substack{S\supseteq T\\|S|\ge2}}|c_j(S)|\,\beta_\chi^{|S|-|T|}.
\]
Expand each character into its allowed-law mean and centred part. Different nonempty
exact leaf supports $T$ are orthogonal in the product reference; one must first sum the
contributions of all centres on the same support, and only then square. For two
distinct centre values $y,w$ outside the middle set, the one-variable character bound
gives
\[
 |\E\phi_y(X)\phi_w(X)|\le\nu_\chi:=q_\varpi+\beta_\chi^2,\qquad
 q_\varpi=\frac{\sum_{s=1}^\varpi\binom{\varpi}{s}(s+\one_{s\ \rm odd})r
 +2^\varpi(z+1)r^2}{1-e_\varpi},
 \tag{PR3}
\]
where $\phi_y(X)=\chi(X-y)-\E\chi(X-y)$; on the diagonal the variance is at most one.
For $s$ independent leaf coordinates and a centre coefficient of modulus at most
$\bar\alpha$, expanding the squared norm over two centre values bounds it by
$\bar\alpha^2(\kappa r^2+\nu_\chi^s)$: the diagonal probability is at most $\kappa/p$,
and off the diagonal the $s$ covariance bounds multiply. For
$s\le2$ the Paley bound $\bar\alpha\kappa r$ is stronger. So
the higher-star marginal, centred at the opposite (input) leaves
rather than at the centre, is at most
\[
 \Bigl[\sum_{T\ne\emptyset}\Bigl(\sum_jb_j(T)\Bigr)^2\mathsf B_{|T|}^2\Bigr]^{1/2},\qquad
 \mathsf B_s^2=\begin{cases}\kappa^2r^2,&s\le2,\\ \kappa r^2+\nu_\chi^s,&s\ge3,\end{cases}
 \tag{PR4}
\]
and the output-centred orientation contributes its empty-input support,
$\sum_j\sum_{|S|\ge2}|c_j(S)|\beta_\chi^{|S|}$. These are constant-image bounds; no
improved full-operator bound is asserted. In the five geometries other than
$({\rm L},2)$ the singleton part of the difference is bounded as the singleton part of
$\mathsf M^{\rm st}_k$, with $\sum_{i'}|M_{\rm F}(i,i')-M_{\rm E}(i,i')|$ in place of
$d_i$.

\paragraph{The six-density Gram matrix.}
For $({\rm L},2)$ the two layer copies of a centre have the same four-pin allowed law.
Index the six distinct densities by $u\in\Zsix$, the object of the other coordinate,
write $\rho_u$ for the density of the allowed set $\mathcal A_u$ relative to the
uniform law and $\psi_u=\rho_u-\one$. For class one, two distinct allowed sets of the
same parity have six compatible pins and opposite-parity sets eight disjoint pins; for
class three, distinct same-parity sets are disjoint, since one common pin has opposed
signs. The leading density Gram matrix $\mathcal H$ has the entries
\[
 \mathcal H_{uu}=16,\qquad \mathcal H_{uu'}=\begin{cases}
 4,&u\ne u',\ u\equiv u'\pmod2,\ p\equiv1\pmod4,\\
 0,&u\ne u',\ u\equiv u'\pmod2,\ p\equiv3\pmod4,\\
 1,&u\not\equiv u'\pmod2,\end{cases}
\]
with absolute entry errors at most
\[
 \eta_{\rm diag}=\frac{16e_4}{1-e_4},\quad
 \eta_{\rm same}=\frac{4(e_6+2e_4+e_4^2)}{(1-e_4)^2},\quad
 \eta_{\rm cross}=\frac{e_8+2e_4+e_4^2}{(1-e_4)^2};\tag{PR5}
\]
the same-parity class-three error is zero. The Gram matrix of the $\psi_u$ is
$\mathcal G=\mathcal H-\one\one^{\mathsf T}$ with the same errors $\eta_{uu'}$. For the
singleton marginal, collapse the twin layer coefficients in each output row to a vector
$\mathbf c_i\in\R^6$. The Paley operator annihilates constants and has norm $r$;
restricting its output to an allowed set costs $\sqrt\kappa$, and distinct centred
output coordinates are orthogonal. Hence the squared singleton marginal is at most
\[
 \kappa r^2\sum_i\Bigl[\mathbf c_i^{\mathsf T}\mathcal G\mathbf c_i
 +\sum_{u,u'}|c_{iu}c_{iu'}|\eta_{uu'}\Bigr].\tag{PR6}
\]
The exact leading quadratic sums are $2448$ and $2064$ in the two prime classes; the
certificate checks both matrix orientations. One may also bound the singleton
\emph{mean} by $r\operatorname{tr}(|\mathbf D|\mathcal G)$: in class one the collapsed
symmetric matrix has first row $(0,4,-2,0,-2,4)$ and absolute-value first row
$(16/3,-2/3,4/3,-8/3,4/3,-2/3)$; in class three its Hermitian absolute value is
$2\sqrt3$ times the projection onto the parity-wise mean-zero vectors. These give the
leading constants $528$ and $128\sqrt3$. The stronger complete mean (PR2) is the one
used for the normalizer.

\paragraph{Weighted moving-root moments.}
For a set $S$ of two to five distinct field elements let $\tau_S$ be the moving-root
trace \eqref{eq:moving-root-tau}, and let $\chi_A$ be a squarefree character weight
with at most four roots. By \cref{cor:moving-root-uniform}, with arbitrary overlaps of
the three root sets,
\[
 \Bigl|\E_x\chi_A(x)\tau_S(x)\tau_{S'}(x)-\frac{\one_{S=S',\ A=\emptyset}}p\Bigr|
 \le512p^{-3/2} ,\tag{PR7}
\]
and, by \cref{cor:moving-root-conditional}, for an allowed law with at most four pins
and density normalization $\gamma$, and distinct root sets $S$,
\[
 \Bigl\|\sum_S\alpha_S\tau_S\Bigr\|_{\rm allowed}\le\sqrt\gamma\,r
 \Bigl[\sum_S\alpha_S^2+8192r\Bigl(\sum_S|\alpha_S|\Bigr)^2\Bigr]^{1/2}. \tag{PR8}
\]
The input is the local monodromy of the moving-root sheaves at their finite roots
(Katz--Sarnak, Remark~10.1.17, \cref{lem:dl:sheaf}): distinct root sets give sheaves
that are not isomorphic after any rank-one twist.

\paragraph{Signed grouping of the principal supports.}
For each exact support $T$ separate its coefficient with no complementary leaf means,
$\alpha_j=\prod_{i\in T}M_{\rm F}(i,j)-\prod_{i\in T}M_{\rm E}(i,j)$; all terms containing
complementary means keep their majorants from (PR4), and no cancellation of unknown
means is assumed. For $|T|\ge3$ put $\Psi=\sum_j\alpha_j\rho_j$, with $\rho_j$ the
density of the allowed law of the centre $j$. The squared norm of its
principal term is exactly $\E_{y,w}\Psi(y)\Psi(w)\prod_{i\in T}\E\phi_y(X_i)\phi_w(X_i)$,
and the diagonal/off-diagonal split with (PR3) bounds it by
\[
 r^2\norm\Psi_2^2+\nu_\chi^{|T|}\Bigl(\sum_j|\alpha_j|\Bigr)^2;\tag{PR9}
\]
after the twin centres are collapsed into a coefficient vector $\mathbf c\in\R^6$,
(PR5) bounds $\norm\Psi_2^2$ by
$\mathbf c^{\mathsf T}\mathcal H\mathbf c+\sum|c_uc_{u'}|\eta_{uu'}$. For a pair support
first replace every centre normalization by one; since
$|1-1/\gamma_j|\le e_4$ for the density normalization
$\gamma_j$ of the allowed law of the centre $j$
(\cref{cor:moving-root-conditional}), the error is at most
$e_4\kappa r\sum_j|\alpha_j|$. Group identical nonempty pin subsets by
$c_S=\sum_{j:S\subseteq{\rm pins}(j)}\alpha_j\prod_{s\in S}{\rm sign}(j,s)$. Fixing the
first leaf makes the moving trace root set $S\cup\{x\}$, whose roots are distinct because
leaves avoid every middle label; (PR8) bounds the nonconstant principal pair term by
\[
 \sqrt{\gamma_0}\,r\bigl[\mathfrak a_T+8192r\,\mathfrak b_T\bigr]^{1/2},\qquad
 \mathfrak a_T=\sum_Sc_S^2,\quad\mathfrak b_T=\Bigl(\sum_S|c_S|\Bigr)^2,\quad
 \gamma_0=(1-e_4)^{-1}.\tag{PR10}
\]
Centring both leaves contracts this bound. Deleted centre labels cost at most
$12\cdot16r^2\sum_j|\alpha_j|$. The empty pin monomial is $\one_{x=y}-1/p$; after both
leaf centrings and restriction to injective output tuples its remaining absolute value is
at most $3\kappa r^2|\sum_j\alpha_j|$, a remainder charged by the triangle inequality
\emph{after} restriction, not an unrestricted product-space norm estimate. For singleton
supports group the pin subsets of size two to four in the same way, with the
coefficients of $\mathbf D$; put $\mathfrak a=\sum_{i,S}c_{iS}^2$,
$\mathfrak b=\sum_i(\sum_S|c_{iS}|)^2$ and $U=\sum_i(\sum_j|\mathbf D_{ij}|)^2=432$. The
principal singleton norm is at most $\sqrt{\gamma_0}\,r\sqrt{\mathfrak a+8192r\mathfrak b}$,
and the centre-normalization and deleted-label errors are at most
\[
 \frac{e_4}{1-e_4}\sqrt{U\gamma_0}\,r\sqrt{11+8192\cdot121r},\qquad
 \gamma_0\,12\cdot16\sqrt U\,r^2.\tag{PR11}
\]
One-pin monomials are constant on the leaf allowed set and disappear under centring.
The direct singleton terms and the size-one components of higher stars lie in the
singleton subspace, which is orthogonal to all supports of size at least two: sum their
squared bounds, take one square root, then add the pointwise pair-diagonal and scalar
empty-support errors. The exact pair-support totals $(\sum_T\mathfrak a_T,\sum_T\mathfrak
b_T)$ are $(6984,221400)$ and $(6792,213144)$ in the two prime classes, and the leading
squared higher-support Gram totals are $18048$ and $18240$; both orientations are
enumerated.

\paragraph{Strict injection, normalization and the differences.}
By \cref{lem:restriction}(c), with the full star bounds $\mathsf C_{\rm F},\mathsf
C_{\rm E}$ and the nonstar and coincidence remainders, a product mean-difference bound
$\Delta_0$ and a restricted-output centred marginal bound $\mathsf M_0$ (from
(PR4)--(PR11)) give
\begin{equation}\tag{PR12}
\begin{aligned}
 |Z_{\rm F}-Z_{\rm E}|\le\Delta&:=\frac{\Delta_0+(\mathsf C_{\rm F}+\mathsf C_{\rm E})
 (2\sqrt{1-I_0}+1-I_0)}{I_0}+\mathsf R,\\
 \norm{P_0(K_{\rm F}-K_{\rm E})\one}\le\mathsf M_\Delta&:=\frac{\mathsf M_0
 +(\mathsf C_{\rm F}+\mathsf C_{\rm E})\sqrt{1-I_0}}{I_0}+\mathsf R,
\end{aligned}
\end{equation}
with $\mathsf R=N^{\rm ns}_{\rm F}+D^{\rm col}_{\rm F}+N^{\rm ns}_{\rm E}+D^{\rm col}_{\rm E}$.
Here $\Delta_0=\norm{\mathbf D}_1\mathsf J_\varpi r^2+\mathsf E_*$, with $\norm{\mathbf D}_1$ the
absolute entry sum of the signed singleton difference ($72$ for $({\rm L},2)$) and
$\mathsf E_*=\sum_j\sum_{|S|\ge2}|c_j(S)|\beta_\chi^{|S|}$ summed over both orientations:
the cancellation in the difference is used before the triangle inequality. Restriction
contracts the orthogonal principal sum, and the pair-diagonal remainder has already been
bounded on the restriction. By \cref{lem:recentering}(c),
\[
 \norm{\varrho_{\rm F}-\varrho_{\rm E}}_2\le\frac{\mathsf M_\Delta
 +\Delta\,\mathsf M_{\rm E}/(1-\eps_{\rm E})}{1-\eps_{\rm F}} .\tag{PR13}
\]
If $\varrho_{\rm E}\ge e^{-\ell_{\rm E}}$, where
$\ell_{\rm E}=\ell_1+\ell_2+\ell_3$ with $\ell_1$ the
conditional completion lower error (the logarithmic deficit in
the lower bound for the conditional completion density),
$\ell_2$ the half lower error (at most
$-n\log(1-e_\varpi)-\log I_0$), and $\ell_3$ the
full-given-middle upper error (the logarithmic excess in the
upper bound for the full density given the middle row),
multiplication by $e^{\ell_{\rm E}/2}$ turns (PR13) into
the relative chi-square bound $\norm{\varrho_{\rm F}/\varrho_{\rm E}-\one}_{L^2(\varrho_{\rm
E}\mu)}\le c_{\rm mar}r$; this defines the \emph{marginal coefficient} $c_{\rm mar}$. Both
conditional normalizer ratios are at most $1+\Delta/(1-\max(\eps_{\rm F},\eps_{\rm E}))$,
and the global normalizer is an average under the common middle reference; hence the
\emph{middle log ratio}
\[
 \ell_{\rm mid}=\frac{2\Delta}{1-\max(\eps_{\rm F},\eps_{\rm E})}
\]
bounds both middle-density log ratios. This linear majorant, rather than a pointwise
logarithmic formula, has the required monotonicity: every coefficient above, after its
indicated power of $r$ is removed, is admissible (\cref{lem:majorant}). Indeed each is
built by sums, products, maxima, minima, square roots, exponentials and the maps
$y\mapsto1/(1-y)$ and $y\mapsto-\log(1-y)$ from fixed integers, $r$ and the admissible
functions $e_t$, $\kappa$, $\beta_\chi/r$, $q_\varpi/r$, $\mathsf T_d$, $\mathsf E_d/r^2$,
$\sqrt{1-I_0}/r$, $N^{\rm ns}_k/r^2$ and $D^{\rm col}_k/r^2$
(\cref{prop:pinned-extraction,prop:pinned-collision}) and the count logarithms in
$\ell_{\rm E}$; so $\eps_k/r$, $\mathsf M_k/r$, $\Delta/r^2$, $\mathsf M_\Delta/r$,
$\ell_{\rm mid}/r^2$ and $c_{\rm mar}$ are admissible. For example
$\mathsf B_s/r=(\kappa+(\nu_\chi/r)^sr^{s-2})^{1/2}$ for $s\ge3$.

\subsubsection{Exact-support correlations}\label{subsec:row-anova-correlation}
The fold row also needs the correlation $\rho_{\rm F}$ of the conditional outer rows in
their own marginals, and the matching row the correlation $\rho_{\rm E}$. Decompose the
mean-zero product space as $\mathcal V\oplus\mathcal W$, where $\mathcal V$ is the
orthogonal sum of the single-coordinate spaces and $\mathcal W$ consists of the supports
of size at least two. For a left centre $i$ write $N_i$ for its neighbours. Summing all
higher stars with exact leaf support $T$ gives an operator of norm at most
$\kappa r(1+\beta_\chi)^{|N_i|-|T|}$ for $|T|\ge2$, and for a singleton leaf the
higher-star bound is $\kappa r((1+\beta_\chi)^{|N_i|-1}-1)$: retain one centred Paley
factor and use the independent remaining leaves as positive Gram factors with diagonal at
most one; complementary leaf means are bounded multipliers on the centre coordinate, and
projecting the centre to mean zero contracts the norm. Let $\mathsf d$ be the sum of the
scalar matrix bounds of the single-edge adjacency matrix and of the two singleton-support
higher-star matrices, each bounded by the square root of its maximal row sum times its
maximal column sum. For the higher-support block put
\[
 \mathsf G_{ii'}=\sum_{s=2}^{|N_i\cap N_{i'}|}\binom{|N_i\cap N_{i'}|}s
 (1+\beta_\chi)^{|N_i|+|N_{i'}|-2s},\qquad
 \mathsf a=\Bigl(\max_i\sum_{i'}\mathsf G_{ii'}\Bigr)^{1/2},
\]
and $\mathsf b$ likewise with the transposed neighbourhoods. Orthogonality of the centre
coordinates and of the leaf supports bounds the two off-diagonal blocks by $\kappa r\mathsf a$
and $\kappa r\mathsf b$, and empty supports vanish after global centring. So the full
centred star operator has norm at most
\begin{equation}\label{eq:row-anova-stars}
 C_\star=\frac{\kappa r}2\Bigl(\mathsf d+\sqrt{\mathsf d^2+4\max(\mathsf a,\mathsf b)^2}\Bigr).
\end{equation}
No sign or symmetry assumption is needed for this scalar block majorant. For the largest
$\rm E$ row, $\mathsf d\to4$ and $\mathsf a^2=\mathsf b^2\to26$, giving the leading
coefficient $8(4+\sqrt{120})<119.636$; for the largest $\rm F$ row the values are $6$ and
$162$, giving $8(6+\sqrt{684})<257.228$. A centred test on the strict reference, extended
by zero with the normalizing square-root factor, remains centred on $\lambda$
(\cref{lem:restriction}(a),(b)); hence the compressed raw kernel $K_k$ satisfies
$\norm{P_0K_kP_0}\le c_{\rm st}:=C_\star/I_0+N^{\rm ns}_k+D^{\rm col}_k$. With
$|Z_k-1|\le\eps_k$ and $\norm{P_0K_k\one},\norm{P_0K_k^*\one}\le\mathsf M_k$ from
\cref{subsec:push:layer72product}, and actual marginal densities at least
$e^{-\ell_{\rm low}}$, \cref{lem:recentering}(b) bounds the conditional correlation by
\begin{equation}\label{eq:row-anova-actual}
 e^{\ell_{\rm low}}\Bigl(\frac{c_{\rm st}}{1-\eps_k}+\frac{\mathsf M_k^2}{(1-\eps_k)^2}\Bigr),
\end{equation}
where $\ell_{\rm low}$ is the insertion-count lower-density bound of the row. The minimum
with the pinned-influence bound \ref{it:record:upper-F2} is taken. These bounds apply to
$\rm F$ and to $\rm E$, and conditional correlations tensorize by taking their maximum
over the primes of a block (\cref{lem:row-comparison}(c)). The scalar majorant
\eqref{eq:row-anova-stars} uses no property of the prime class; it is the star bound of
the tail datum.

\paragraph{Full star matrices (bands).} At $p\equiv3\pmod4$ the centred star part of a
pinned row is bounded by the full star matrix under the pinned product law
(\cref{prop:conditional-star}): the positive semidefinite gate
$\begin{psmallmatrix}\mathsf k^2I-AA^*&\mathsf kD\\\mathsf kD^*&\mathsf k^2I-BB^*\end{psmallmatrix}
\succeq0$ gives the norm bound $\kappa_0r(\mathsf k+\mathsf e_\star)$, where $\kappa_0\ge\kappa$
is the assigned rational density bound and the decay weights $\omega_s$ of
\cref{prop:conditional-star}(c) multiply the columns of the supports of size $s$. It is
transported to the actual marginals by \eqref{eq:row-anova-actual}, with
$\min\{C_\star,\kappa_0r(\mathsf k+\mathsf e_\star)\}$ in place of $C_\star$. The bands use
this form in \ref{it:record:upper-F4}; the tail datum uses
\eqref{eq:row-anova-stars} in \eqref{eq:row-anova-actual}, which hold at both classes.
The gates are rational and do not depend on $p$, and the other normalized errors are
admissible, so a table certified at an endpoint holds at every larger class-three prime
(\cref{cert:upper-band-inputs}(c)).

\subsubsection{Fold and matching conditional laws bounded jointly}\label{subsec:upper-joint-rows}
For the bands, the fold and matching rows of a coordinate are bounded jointly, on one
common pinned product reference transported to the actual $\rm E$ marginal
(\cref{lem:joint-rows,cor:joint-row-split}): one Gram matrix bounds the fold kernel and
the matching kernel together, and the completion of \cref{lem:five-type}(f) spends the
reserve of the matchings on that matrix rather than on the sum of two separate norms.
For the $\rm L$ rows the pin classes are mutually exclusive, and this lowers the
pure-star coefficient.

\begin{lemma}[Pin classes of the $\rm L$ rows]\label{lem:joint-pins}
Let $p\equiv3\pmod4$ and consider the conditional $\rm L$ rows with $\ell\in\{1,2\}$
layers (the geometries $({\rm L},\ell)$ of \cref{subsec:upper-rows}). Index the
labels of a row by the six base coordinates $u\in\Zsix$ and, for $\ell=2$, their two
\emph{twins} (one per layer). The one-layer signed singleton matrices of the fold and
the matching law are
\[
 D_{\rm F}=\begin{pmatrix}
 0&1&1&0&0&1\\1&0&1&0&0&1\\0&1&0&1&1&0\\
 0&1&1&0&1&0\\1&0&0&1&0&1\\1&0&0&1&1&0\end{pmatrix},\qquad
 D_{\rm E}=\begin{pmatrix}
 0&0&1&0&-1&0\\0&0&0&-1&0&1\\-1&0&0&0&1&0\\
 0&1&0&0&0&-1\\1&0&-1&0&0&0\\0&-1&0&1&0&0\end{pmatrix};
\]
for $\ell$ layers tensor each with the $\ell\times\ell$ all-ones matrix $J_\ell$, and
form $\mathcal Z$, $\mathcal A$ and $\mathcal G=\mathcal A\mathcal A^*$ as in
\cref{lem:joint-rows}, with $\mathcal Z_0$ the matrix $\mathcal Z$ for $\ell=1$. Let $\Pi$
be the projection onto the vectors that are constant on each twin class, in both copies
of the joint block.
\begin{enumerate}[label=\textup{(\alph*)}]
\item $|\mathcal Z|=(-\mathcal Z^2+3\ell^2\Pi)/(4\ell)$; in particular
 $|\mathcal Z|\pm i\mathcal Z\succeq0$.
\item Let $\mathcal S$ consist of the nine coordinate sets obtained by choosing one
 even and one odd base coordinate, with all their twins in both copies. At every label
 the coordinates whose allowed set contains it lie in one set of $\mathcal S$.
\end{enumerate}
Consequently, let $T$ be $|\mathcal Z|$, or a real symmetric matrix with
$T\pm i\mathcal Z\succeq0$ (for $\ell=2$ it suffices that $T=T_0\otimes J_2$ with
$T_0\pm i\mathcal Z_0\succeq0$). If $\mathsf k>0$ and the rational gates
$\mathsf k^2I_S-\mathsf k\kappa_0T_S-\kappa_0^2\mathcal G_S\succeq0$ hold for all
$S\in\mathcal S$, then $r\mathsf k$ is a joint pure-star coefficient: the lower bound of
\cref{lem:joint-rows}(a) holds with $r\mathsf k/I_0$ in place of $\kappa_0r\mathsf k/I_0$
in $\mathfrak k_p$, all other error and normalizer charges being kept.
\end{lemma}

\begin{proof}
\emph{The matrices.} Between the outer objects $0$ and $2$ the fold law $\rm F$ has the
forward pair $0\to2$ and no relation $2\to0$, so its cross edges are $(0,u)\to(2,v)$
for $v\in\{u+1,u-1,\rot u\}$; the matching law $\rm E$ has both relations and no
forward pair, so its cross edges are $(0,u)\to(2,\rot u)$ and $(2,v)\to(0,\rot v)$
(\cref{subsec:push:layer72spectral}). With the sign $+1$ for an edge from the row $0$
to the row $2$ and $-1$ for the reverse, these are the displayed $D_{\rm F}$ and
$D_{\rm E}$; edges carry no layer condition, which gives the factor $J_\ell$.

(b) An edge between an outer object ($0$ or $2$) and the middle object $1$ of the
chamber coordinate needs a forward relation in the other coordinate, which has type
$\rm L$, and $x\to y$ is forward there only for $y=\rot x$
(\cref{subsec:push:layer72spectral}). So in both laws (\cref{lem:upper-geometry}(d))
the outer label of base coordinate $u$ is joined to the middle labels $m_v$ exactly by
the edges $u\to m_{\rot u}$ and $m_{\rot^{-1}u}\to u$. For the six base coordinates
the out-pin and the in-pin are
\[
 (0)\ m_2,m_4;\quad(1)\ m_5,m_3;\quad(2)\ m_4,m_0;\quad
 (3)\ m_1,m_5;\quad(4)\ m_0,m_2;\quad(5)\ m_3,m_1 .
\]
Two even base coordinates share a pin with opposite orientations ($0$ and $2$ at
$m_4$, $0$ and $4$ at $m_2$, $2$ and $4$ at $m_0$), and so do two odd ones. Since
$p\equiv3\pmod4$, the strict-square relation is a tournament, and no label $t$
satisfies both $t\to m$ and $m\to t$; so at every label at most one even and one odd
base coordinate is admissible. With two layers every edge between layers is present
(there is no condition on the layers), so the two twins of a base coordinate have the
same pins and the same allowed set. The two outer rows have the same pins
(\cref{lem:upper-geometry}(d)), so both copies of the joint block have the same
allowed sets. Excluding the middle labels and the other fixed labels only shrinks the
allowed sets. This proves (b).

(a) For $\ell=1$, $\mathcal Z$ is a real antisymmetric $12\times12$ matrix, and direct
multiplication of the displayed matrices gives $(\mathcal Z^2+I)(\mathcal Z^2+9I)=0$.
So the eigenvalues of $-\mathcal Z^2$ are $1$ and $9$, those of
$|\mathcal Z|=(-\mathcal Z^2)^{1/2}$ are $1$ and $3$, and $|\mathcal Z|=(-\mathcal Z^2+3I)/4$.
For $\ell=2$, $\mathcal Z=\mathcal Z_0\otimes J_2$ with $J_2^2=2J_2$ and $\Pi=I\otimes J_2/2$,
so $|\mathcal Z|=|\mathcal Z_0|\otimes J_2=\frac14(-\mathcal Z_0^2+3I)\otimes J_2
=(-\mathcal Z^2+12\Pi)/8$. As $\mathcal Z$ is normal and $i\mathcal Z$ is Hermitian,
$|\mathcal Z|\pm i\mathcal Z\succeq0$ on every eigenvector of $\mathcal Z$. If
$T_0\pm i\mathcal Z_0\succeq0$, then $(T_0\pm i\mathcal Z_0)\otimes J_2\succeq0$, that is,
$T\pm i\mathcal Z\succeq0$ for $T=T_0\otimes J_2$.

For the consequence, a coordinate $(c,i)$ of $\mathcal Z$ has $\rho_i(t)>0$ exactly when
its allowed set contains $t$; so by (b), for every label $t$ these coordinates lie in
one member of $\mathcal S$, which is the hypothesis on $\mathcal S$ in
\cref{lem:joint-potential}, and by (a) the matrix $T$ satisfies its hypothesis
$T\pm i\mathcal Z\succeq0$. So \cref{lem:joint-potential} applies, and its conclusion is
the assertion. The Schur elimination of the higher features and the pointwise
diagonal contraction on the active coordinates are made there.
\end{proof}

The tables of \cref{cert:upper-band-inputs}(c) record, for each of the six row
geometries, the joint coefficient $k_{\rm J}$, the common fold coefficient $c_{\rm F}$
and the matching coefficient of \cref{lem:joint-rows}, and, for the $\rm L$ rows, a
rational base potential $T_0$ with the nine gates of \cref{lem:joint-pins}.

\begin{proposition}[A necessary floor for the chosen joint-row formula]
\label{prop:joint-row-floor}
Every admissible two-layer $\rm L$ table above $2^{47}$, using the
potential and decay gates of \cref{lem:joint-pins,prop:conditional-star},
has chosen complete coefficient $k_{\rm J}>172.9$.
Consequently the band completion inputs satisfy
$f_{\rm L}+h_{\rm L}M_{\rm L}>172.9$, and the native maximum
recursion for the reference gives
$w_{\rm LR},w_{\rm RL}>172.9d_{\rm L}$.
These are lower bounds for specified certificate formulas and their
chosen weights, not for the underlying operator norms. The tail
uses different row formulas and is excluded.
\end{proposition}
\begin{proof}
A free coordinate is uniform on a four-pin allowed set of size $n$,
excluding twelve fixed middle labels. The character expansion gives
\[
 \frac{p-11\sqrt p-6}{16}-12\le n\le
 \frac{p+11\sqrt p+6}{16}.
\]
For $p\ge2^{47}$ this implies $p/17\le n\le c_0p$ with
$c_0=(1+11\cdot2^{-23})/16$. The positive $6$ in the upper count
is absorbed by
$11\cdot10^{-7}+6\cdot2^{-47}<11\cdot2^{-23}$.
For $X_t(x)=\chi(t-x)$, Paley orthogonality gives
\[
 \sum_t\E X_t^2=p-1,\qquad
 \sum_t(\E X_t)^2=p/n-1.
\]
After deleting the twelve fixed labels, the covariance matrix has
trace at least $p-29$, rank at most $n-1$, and diagonal at most $1$.
Thus an off-diagonal covariance bound $\nu_\chi$ must satisfy
\[
 p\nu_\chi^2\ge(1-29/p)^2/c_0-1>14.9999.
\]
Put $c=15.9999$. The decay conditions of
\cref{prop:conditional-star}(c) now force
$\kappa_0>c$, $\kappa_0^2\omega_2^2\ge c^2$ and
$\kappa_0^2\omega_s^2\ge c$ for $s\ge3$; a weight chosen as $1$
already satisfies the corresponding inequality.

Let $G_2,G_h$ be the unweighted column Grams in the joint higher-star
matrix, for supports of size $2$ and at least $3$ respectively.
For each of the nine pin classes choose a real vector $v_i$ supported
on it, and let $R$ have these vectors as columns. Put
$M=R^T\mathcal ZR$. If $K$ is a real skew contraction, then
$R^TTR\pm iM\succeq0$ and $I\pm iK\succeq0$ imply
\[
 \tr(R^TTR)\ge\langle K,M\rangle_F.
\]
Indeed take traces of
$(I+iK)(R^TTR-iM)/2$ and $(I-iK)(R^TTR+iM)/2$.
Evaluating the nine potential gates on their vectors and adding
therefore forces
\[
 k^2N-kcB-c^2Q_2-cQ_h\ge0,
\]
where $N=\sum_i\|v_i\|^2$, $Q_j=\sum_i v_i^TG_jv_i$,
and $B=\langle K,M\rangle_F$. The exact rational witness of
\cref{cert:astra-local}, reconstructed from the displayed
$D_{\rm F},D_{\rm E}$, has
\[
 N=10000010408,\quad Q_2=558217933120,\quad Q_h=1038757368976,
 \quad B=\frac{315100239978216987}{6250000}.
\]
At $k=172.9$ the polynomial equals
$-29790250369036028509477/625000000000<0$.
It is negative at zero and convex, hence negative throughout
$[0,172.9]$. The complete coefficient adds nonnegative mean,
remainder and normalization costs to $k$, so $k_{\rm J}>172.9$.
Finally \ref{it:record:upper-F5}(a) gives
$f_{\rm L}\ge H(k_{\rm J}-h_{\rm L})_+$ and
$M_{\rm L}\ge H\ge1$, whence
$f_{\rm L}+h_{\rm L}M_{\rm L}\ge k_{\rm J}$.
Apply the ordinary completion inequalities on both axes.
\end{proof}

\subsection{Band input formulas (F1)--(F6) and (F10): definitions}\label{subsec:upper-formulas}

The band inputs are computed by the following formulas, each stated for every odd
prime at which its listed hypotheses (the \emph{gates}) hold; the formulas
\ref{it:record:upper-F7}--\ref{it:record:upper-F9} of the comparison laws follow in
\cref{sec:upper-comparison}, and \cref{lem:record:domain} extends gates checked at an
evaluation point to every larger class-three prime and to every block of the band. For
a block of primes at least $p_0$ and logarithmic length at most $L'\le10000$ put
$\bar n=L'/\log p_0$.
\begin{enumerate}[label=\textup{(F\arabic*)}]
\item\label{it:record:upper-F1} \emph{Insertion counts} (\cref{lem:extension}). An
 insertion with $t$ signed constraints at distinct labelled roots and $z$ excluded
 labels has normalized count in $[1-e^-_t(z),1+e^+_t]$; Weil's bound is used in both
 parities. An explicit order of the free vertices gives, with $t$ and $z$ the numbers of
 labelled neighbours and of labelled vertices at each step, the logarithmic count
 bounds $c^+=\sum e^+_t$, an upper bound for $\sum\log(1+e^+_t)$, and
 $c^-=\sum-\log(1-e^-_t(z))$. \emph{Hypothesis:} $e^-_t(z)<1$ at every insertion.
\item\label{it:record:upper-F2} \emph{Conditional influence}
 (\cref{cor:influence-schur}, the certified form of \cref{prop:middle-negative}). For
 disjoint free sets $X,Y$ the conditional maximal correlation is at most
 $r(c_\times+\sqrt{|X||Y|}\,rq_{\rm inf}^2/(1-rq_{\rm inf}))$, with the count errors
 $w=e^-_\Delta(n)$ and $w^{\rm conn}=e^-_{\Delta+1}(2n)$, where $n$ is the number of
 vertices of the graph and $\Delta$ the largest degree of a free vertex. \emph{Hypotheses:}
 $w<\frac12$, $w^{\rm conn}<1$, $rq_{\rm inf}<1$, and the resulting bound below one.
\item\label{it:record:upper-F3} \emph{Product rows} (\cref{subsec:push:layer72product,subsec:advance:product-rows}),
 with the pinned nonstar extraction: the normalizer errors $\eps_{\rm F},\eps_{\rm E}$,
 the differences $\Delta$, $\mathsf M_\Delta$, the middle log ratio $\ell_{\rm mid}$ and
 the marginal coefficient $c_{\rm mar}$. A constant that depends on the class of $p$
 is taken as the larger of its values in the two classes
 (\cref{subsec:advance:product-rows}). \emph{Hypotheses:} $e_\varpi<\frac12$ for the pin
 count and $e^{\rm pin}_\varpi<\frac12$ for the pinned extraction
 (\cref{prop:pinned-extraction});
 $I_0=1-\binom n2\kappa r^2>\frac12$; $\eps_{\rm F},\eps_{\rm E}<\frac12$; every nonstar
 denominator positive.
\item\label{it:record:upper-F4} \emph{Exact-support correlation}
 (\cref{subsec:row-anova-correlation}, \eqref{eq:row-anova-actual}), minimized with
 \ref{it:record:upper-F2}: for the bands with the full star matrix of
 \cref{prop:conditional-star}, that is, with $\min\{C_\star,\kappa_0r(\mathsf k+\mathsf e_\star)\}$
 in \eqref{eq:row-anova-actual}, and for the tail datum with \eqref{eq:row-anova-stars};
 class-dependent constants are taken as in \ref{it:record:upper-F3}.
 \emph{Hypotheses:} $1-\eps>0$ for the normalizer error, the selected correlation below
 one, and, for the bands, the gates of \cref{prop:conditional-star}.
\item\label{it:record:upper-F5} \emph{Completion rows.} With $\ell_{\rm mid}$ and
 $c_{\rm mar}$ from \ref{it:record:upper-F3}, all at $p_0$, put
 \[
  H=\exp(\bar n\ell_{\rm mid}),\qquad \upsilon^2=\exp(\bar nc_{\rm mar}^2/p_0)-1 .
 \]
 \emph{(a) Joint rows (bands).} Let $(k_{\rm J},c_{\rm F})$ be the joint and the common
 fold coefficient of one complete joint-row table (\cref{lem:joint-rows},
 \cref{lem:joint-pins}) whose endpoint is at most $p_0$, let $h$ be a proved
 coefficient of the $\rm E$ row (the centred matching kernel has norm at most $hr_p$ at
 every prime $p\ge p_0$, so that the matching block of a block with least prime $b$ is
 at least $-hb^{-1/2}I$ by \cref{lem:row-comparison}(c); for instance
 $h=\rho_{\rm E}/r_0$ of \ref{it:record:upper-F4}), and put
 $E_{\rm T}=(2\upsilon+\upsilon^2)(\exp(\bar nc_{\rm F}r_0)-1)$. Then for every rational
 $t_{\rm Y}>0$ with $t_{\rm Y}H<1$ the coefficient $h$, every $M\ge H+t_{\rm Y}$ and
 every
 \[
  f\ge\frac{H\{(k_{\rm J}-h)_+r_0+E_{\rm T}\}+(H+H^2/t_{\rm Y})\upsilon^2}{r_0}
 \]
 are valid row coefficients (\cref{lem:five-type}(f),
 \cref{cor:joint-completion-blocks}). An ordinary input takes $f$ to be the maximum of its one-layer and two-layer values, computed with one common value of $t_{\rm Y}$, and likewise takes $H+t_{\rm Y}$ to be the maximum of its one-layer and two-layer values; a reserved
 input uses the one-layer row. A smaller proved $\rm E$ bound may be used for $h$,
 since all of them concern the same actual law; the pair $(k_{\rm J},c_{\rm F})$ always
 comes from one table.

 \emph{(b) Schur rows (tail datum).} With moreover the pointwise marginal ratio
 $\ell_{\rm mar}$ of the field records (from \ref{it:record:upper-F1} and
 \cref{lem:four-count} with $M$ the middle row and $S$ an outer row; its hypothesis,
 that the $\rm F$ and $\rm E$ graphs induce the same graph on $M\cup S$, is
 \cref{lem:upper-geometry}(d)) and the correlations $\rho_{\rm F},\rho_{\rm E}$ of
 \ref{it:record:upper-F4}, all at $p_0$, put $R_{\rm m}=\exp(\bar n\ell_{\rm mar})$. Then
 \[
  f_{\rm Schur}=\frac{HR_{\rm m}\rho_{\rm F}+(H+50H^2)\upsilon^2}{r_0},\qquad
  h=\frac{\rho_{\rm E}}{r_0}
 \]
 are valid row coefficients (\cref{lem:five-type}(c) with
 \cref{rem:five-type-fold-data} at $t_{\rm Y}=1/50$), and so is the full-norm
 coefficient $f=[H(R_{\rm m}\rho_{\rm F}+2\upsilon+\upsilon^2)+H-1]/r_0$ of
 \cref{lem:row-comparison}(b) with \cref{lem:five-type}(b); the smaller of the two is
 used.

 \emph{Hypotheses:} in (a), $M\ge H+t_{\rm Y}$ and
 $t_{\rm Y}H<1$ for the matching ratio $M=w_{\rm S}/w_{\rm D}$,
 $c_{\rm F}r_0<1$, and the prime gates of the table at its
 endpoint; in (b) the weaker $t_{\rm Y}\le H^{-1}$ suffices, and
 with $M=21/20$ and $t_{\rm Y}=1/50$ both hypotheses follow from
 $H\le103/100$.
\item\label{it:record:upper-F6} \emph{Nonroot spectral caps}
 (\cref{thm:spectral-criterion}, \cref{subsec:upper-nonroot}). For each of the ten
 nonroot cuts, $-grI\preceq\Op\preceq(1+B_{\rm top}/p)I$ with
 $g\ge\max_\pm(\cperp+\tau^{\rm sp}_{\mathfrak c,\pm})/(1-\bar\eps_V)$, the maximum over
 the signs and the classes tested, and $B_{\rm top}=p_0\eps_\varrho(r_0)^2/(1-gr_0)$;
 moreover $\norm{P_0\Op P_0}\le gr$ and
 $\norm{P_0\Op\one}\le\eps_\varrho(r)$ with $p\,\eps_\varrho(r)^2\le p_0\,\eps_\varrho(r_0)^2$,
 so that $\norm{P_0\Op\one}\le B_{\rm top}^{1/2}r$ (\cref{thm:spectral-criterion}(b),(c));
 a block of primes at least $X\ge p_0$
 and length at most $L'$ has negative coefficient at most $d\,b^{-1/2}$ with
 $d=g\exp(B_{\rm top}L'/(X\log X))$. The observed caps of \cref{prop:observed-caps}
 are evaluated in the same way. \emph{Hypotheses:} (S1) of \cref{def:spec:data}
 at $p_0$; $\bar\eps_V<1$; $\eps_Z<\frac12$ (which implies the hypothesis $\eps_Z<1$ of
 \cref{thm:spectral-criterion}); the exact rational matrix inequality
 \eqref{eq:spectral-test} with the recorded rational $\tau^{\rm sp}_{\mathfrak c,\pm}$,
 for the class three and both Fourier signs (bands), or for both classes and both
 signs (tail datum); $gr_0<1$; and $d/\sqrt{p_0}<1$.
\end{enumerate}
\begin{enumerate}[label=\textup{(F\arabic*)},start=10]
\item\label{it:record:upper-F10} \emph{Completion} (\cref{def:record:upper-bands}(d),
 \cref{lem:upper-reference-rp}). The forty two-layer and forty one-layer slice
 inequalities \eqref{eq:upper-slice}, with the recorded $M$, and the twenty-five
 reserved dominations \eqref{eq:independent-fold-domination}, with inputs evaluated
 for block lengths at most $10000$.
\end{enumerate}

\begin{proof}[Proofs and forms]
\ref{it:record:upper-F1} is \cref{lem:extension}(b),(c) at each step of the order, the
graded errors being those of \eqref{eq:graphs:graded}; the evaluators sum $e^+_t$,
which bounds $\log(1+e^+_t)$. \ref{it:record:upper-F2} is \cref{cor:influence-schur},
with the graded count errors of \cref{rem:influence-uses}(1).
\ref{it:record:upper-F3} and \ref{it:record:upper-F4} are
\cref{subsec:upper-rows}. For \ref{it:record:upper-F5}: let the block have least prime
$b$, and put $r_b=b^{-1/2}\le r_0=p_0^{-1/2}$. On a block the densities and chi-square
defects multiply and the correlations are maximized over the primes
(\cref{lem:row-comparison}(c)); the block has at most $\bar n$ primes, so
$\sum_p1/p\le\bar nr_b^2$, and every per-prime quantity, after its indicated power of $r$
is removed, is at most its value at $p_0$ (\cref{lem:record:domain}(i)). So the middle
density lies in $[H^{-1},H]$, the fold marginals have pointwise ratio at most $R_{\rm m}$,
the block correlations are at most $(\rho_{\rm F}/r_0)r_b$ and $(\rho_{\rm E}/r_0)r_b$,
and the relative $L^2$ defect $\upsilon_b$ of the block fold marginals satisfies
$\upsilon_b^2\le\prod_p(1+c_{\rm mar}^2r_p^2)-1\le\exp(c_{\rm mar}^2\bar nr_b^2)-1\le
(\upsilon^2/r_0)\,r_b$, because $x\mapsto(e^{cx^2}-1)/x$ is nondecreasing and
$\exp(c_{\rm mar}^2\bar nr_0^2)-1=\upsilon^2$; in particular $\upsilon_b\le\upsilon$.

(a) is \cref{lem:five-type}(f) through \cref{cor:joint-completion-blocks}, with these
block values: by the table, the prime bounds of \cref{cor:joint-completion-blocks} hold
at every prime of the block with the coefficients $(k_{\rm J},c_{\rm F})$ and the
matching coefficient of the table, the matching block is at least $-hr_bI$, and the
tensor error is at most $(2\upsilon_b+\upsilon_b^2)(e^{\bar nc_{\rm F}r_b}-1)\le
(E_{\rm T}/r_0)\,r_b$, since $x\mapsto(e^{\bar nc_{\rm F}x}-1)/x$ is nondecreasing; so
the coefficient $f_{\rm J}$ of \cref{cor:joint-completion-blocks} at $r=r_b$ is at most
$f$. The Young allowance $t_{\rm Y}$ enters only through $H+t_{\rm Y}$ and
$H^2/t_{\rm Y}$, so it may be chosen freely, row by row, subject to $t_{\rm Y}H<1$.

(b) \Cref{rem:five-type-fold-data}, with these block values and $r=r_b$, turns them
into the fold data of \cref{lem:five-type}(c) at $t_{\rm Y}=1/50$ with
$f_{\rm al}\le f_{\rm Schur}$, and the matching row has exact marginals with
$\norm{T_{\rm E}-\Jproj}\le hr_b$; the full-norm coefficient scales in the same way.
For the tail datum $M=\frac{21}{20}=\frac{103}{100}+\frac1{50}\ge H+t_{\rm Y}$ and
$t_{\rm Y}=\frac1{50}\le\frac{100}{103}\le H^{-1}$ when $H\le103/100$; also $M\ge1$, so
\cref{lem:five-type}(b) applies to the full-norm coefficient. In every row of the tail
datum the Schur coefficient is the smaller one (\cref{rem:five-type-certificates}).

\ref{it:record:upper-F6} is \cref{thm:spectral-criterion} with the standard inputs of
\cref{subsec:upper-nonroot}, for each class tested, and its block form is
\cref{cor:spectral-criterion-blocks} with $P=X$, with $g\ge g_3$ for the bands and
$g\ge\max(g_1,g_3)$ for the tail datum. \ref{it:record:upper-F10} is used in
\cref{lem:upper-reference-rp}.
\end{proof}

For the root cuts the domain of the caps is that of
\cref{lem:record:monotone-endpoint} (route~B) and of \cref{lem:upper-completed-root}
(route~A); the per-prime constants of the nonroot cuts are also read by the one-way
repairs \ref{it:record:upper-F7}.

\subsection{(CP1) for upper atoms}\label{subsec:upper-reference-rp}

We say that the inputs of an upper band are \emph{valid} on a finite set $\Blk$ of
primes, with least prime $b$ and $r=b^{-1/2}$, if for every field type $\varphi$:
\begin{enumerate}[label=(V\arabic*)]
\item\label{it:upper-valid-root} the normalized block cut operator of the chamber cut
 of the $\rm L$ coordinate of $G_{{\rm L}\varphi,2}$ is at least $-d_\varphi r$, and
 that of $G_{{\rm L}\varphi,1}$ at least $-d^{\rm res}_\varphi r$;
\item\label{it:upper-valid-rows} at every value of the middle object, the two-layer
 rows at $\varphi$ satisfy, with graded coefficients at most $(f_\varphi,h_\varphi)$,
 one of three sets of hypotheses: those of \cref{lem:five-type}(f) with
 $(M_\varphi,t_\varphi)$; those of \cref{lem:five-type}(c) with $t_{\rm Y}=t_\varphi$
 and $M_\varphi\ge H+t_\varphi$, $t_\varphi H\le1$; or those of
 \cref{lem:five-type}(b). The one-layer rows satisfy the same with the reserved input;
\item\label{it:upper-valid-layer} for every field pair $\varphi\varphi'$, the normalized
 block cut operator of the layer cut of $G_{\varphi\varphi',2}$ is at least
 $-k_{\varphi\varphi'}r$. For each state $\mathfrak t$ with
 $k_{\mathfrak t}<k_{\mathfrak f(\mathfrak t)}$, the compression of the normalized
 block cut operator of EL22 to the functions of the labels that the components of
 $\mathfrak t$ observe (the merge or the constant labels of \cref{prop:observed-caps})
 is at least $-k_{\mathfrak t}r$.
\end{enumerate}

\begin{lemma}[The upper reference]\label{lem:upper-reference-rp}
\leavevmode
\begin{enumerate}[label=(\roman*)]
\item Let $\Blk$ be a finite set of odd primes on which the inputs of an upper band are
 valid, and let $\mu_\Blk$ be defined by the formula of \cref{def:upper-reference} with
 the weights of the band. Then $\mu_\Blk$ is a probability measure on $X_\Blk^V$; it is reflection
 positive at every cut of $\mathfrak F$; it is invariant under the square-affine maps of
 \ref{it:record:R3}; its one-slot marginals are uniform; and it is admissible
 (\cref{def:admissible}).
\item The inputs of every upper band are valid on every block of the band. Hence every
 upper block satisfies \ref{it:record:R1}--\ref{it:record:R3}, with \ref{it:record:R2} at
 every cut of $\mathfrak F$.
\end{enumerate}
\end{lemma}

\begin{proof}
(i) \emph{Measure, invariance, marginals.} Each component is the image of a field law
$\lambda_\Blk(G)$ under an emission, and $\operatorname{Diag}_\iota$ is an emission
too; the field laws are square-affine invariant at every prime, and images, products
over the primes of $\Blk$ and nonnegative combinations preserve invariance
(\cref{lem:marginals}(1)). So $\widetilde\mu_\Blk$ is a nonnegative invariant measure,
of mass $P_\mu(r)$ since $\pi_{\mathfrak t,j}$ has mass $m_{\mathfrak t}$, and $\mu_\Blk$
is a probability measure with uniform one-slot marginals (\cref{lem:marginals}(2)).

\emph{Chamber cuts.} Fix a chamber coordinate $c$ and a type $\mathfrak t'$ of the other
coordinate. The two-layer row data of $c$ at $\mathfrak f(\mathfrak t')$ form a five-type datum
(\cref{subsec:layer72:upper}), with $\pi_{\rm F}$, $\pi_{\rm E}$ the row laws of the
$\rm F$ and $\rm E$ field laws, $\Gamma$ and $\nu$ the laws of the rows $0,1$ and of the
row $0$ under the $\rm E$ law, and $\nu^b$ its conditional law of the row $0$; the merges
and the constants use only these rows, which is the role of the hidden objects. The
weights of the slice in $\widetilde\mu_\Blk$ are $w_{\mathfrak t}r^{\deg\mathfrak t}=
r^{\deg\mathfrak t'}\,a_{\mathfrak s}r^{\deg\mathfrak s}$ ($\mathfrak t_c=\mathfrak s$, $\mathfrak t_{1-c}=\mathfrak t'$,
$a_{\mathfrak s}=w_{\mathfrak t}$): $r^{\deg\mathfrak t'}$ times the graded weights of
\cref{lem:five-type}(b),(c),(f) formed from the $a_{\mathfrak s}$. The conditions are homogeneous,
so the factor $r^{\deg\mathfrak t'}$ does not matter. By
\ref{it:upper-valid-root} the root envelope $\Pi_{\rm L}\succeq_{\vartheta_0}
-d\,r\,\Pi_{F_0}$ holds with $d=d_{\mathfrak f(\mathfrak t')}$, the half law of the cut being
the law of the aligned fold (\cref{lem:upper-geometry}(c)); by \ref{it:upper-valid-rows}
the rows satisfy the hypotheses of \cref{lem:five-type}(f), (c) or (b) with the input
$(f,h,M,t)$ of $\mathfrak f(\mathfrak t')$; and \eqref{eq:upper-slice} is the graded form of
their conditions: $a_{\rm D}\ge d\,a_{\rm L}$; $a_{\rm S}\ge M\,a_{\rm D}$ with the recorded
$M\ge H+t_{\rm Y}$ and $t_{\rm Y}\le H^{-1}$, so in particular $a_{\rm S}\ge a_{\rm D}$;
$a_{\rm R}\ge fa_{\rm D}+ha_{\rm S}$; and $a_{\rm C}a_{\rm S}\ge a_{\rm R}^2$.
By \cref{lem:five-type}(d), the two-layer sum $\sum_{\mathfrak t}w_{\mathfrak t}
r^{\deg\mathfrak t}\pi_{\mathfrak t,2}$ is reflection positive at every chamber cut of
both coordinates. The same argument with the one-layer row data, the reserved inputs
and the weights $v_{\mathfrak t}r^{\deg\mathfrak t}$ shows that $\sum_{\mathfrak t}
v_{\mathfrak t}r^{\deg\mathfrak t}\pi_{\mathfrak t,1}$ is reflection positive at the chamber
cuts, and $\operatorname{Diag}_\iota$ preserves this (\cref{lem:five-type}(e)). Cut forms
add, so $\widetilde\mu_\Blk$ is reflection positive at every chamber cut.

\emph{The reserved cut.} By \ref{it:upper-valid-layer}, for the field law of each field
pair $\lambda\succeq_\iota-k\,r\operatorname{Diag}_\iota(\lambda^{(1)})$, the half graph of
the layer cut being the one-layer graph (\cref{lem:upper-geometry}(c)); by the last
assertion of \cref{lem:five-type}(e) every two-layer component of $\mathfrak t$ inherits
$\Pi\succeq_\iota-k_{\mathfrak f(\mathfrak t)}r\operatorname{Diag}_\iota(\Pi^{(1)})$, the
merges and constants included. A merge or constant component reads only the observed
rows of its $\rm E$ coordinate, so its cut form at $\iota$ is that of its field law on
the functions of the observed labels; by the second assertion of
\ref{it:upper-valid-layer} the components of the states with
$k_{\mathfrak t}<k_{\mathfrak f(\mathfrak t)}$ satisfy the same order with $k_{\mathfrak t}$,
a cap of their own, which \cref{lem:five-type}(e) allows. With
\eqref{eq:independent-fold-domination}, $v_{\mathfrak t}\ge k_{\mathfrak t}w_{\mathfrak t}$
with the cap of \ref{it:upper-valid-layer} for the component of $\mathfrak t$,
\cref{lem:five-type}(e) gives $\widetilde\mu_\Blk\succeq_\iota0$. Normalization by
$P_\mu(r)>0$ preserves reflection positivity, which is (R2) at every cut of $\mathfrak F$.

\emph{Admissibility.} $\mu_\Blk$ is nonnegative and invariant, with uniform whole-slot
marginals, and it is a law of the field labels alone, independent of the interval,
residue and inclusion indices; this is (O1) of \cref{def:admissible}. The reference
creates no sharp operator bound, index projection or fold, so (O2)--(O4) are those of
the frame.

(ii) Let $\Blk$ be a block of the band $[X,Y)$. Its primes are $3$ modulo $4$, at least
$X$, and $\log q_\Blk\le L_{\rm band}<L_{\rm u}$ (\cref{def:record:upper-bands}(b)). The
nonroot and observed caps were evaluated at $p_0=2^{47}$ or $2^{48}$, not above $X$, and
inflated at $X$ with the length $L_{\rm band}$; the product-row quantities at $p_0=X$
with the length $L_{\rm band}$, and the joint coefficients at an endpoint not above $X$;
all under the gates of \ref{it:record:upper-F1}--\ref{it:record:upper-F6}
(\cref{cert:upper-band-inputs}, \cref{cert:upper-lower-bands}). By
\cref{lem:record:domain}(i) the gates hold at every prime of $\Blk$ and every per-prime
constant is at most its value at $p_0$; the block forms of \ref{it:record:upper-F5},
\ref{it:record:upper-F6} and \cref{prop:observed-caps} (\cref{lem:record:domain}(ii),
with the inflation taken at $X$, since every prime of $\Blk$ is at least $X$) bound the
negative coefficients of the nonroot cuts and of the observed marginals on $\Blk$ by
$d\,b^{-1/2}$, and the row coefficients by $f\,b^{-1/2}$ and $h\,b^{-1/2}$ with the
recorded $M$ and $t_\varphi$; the inputs are upper bounds for these values. Inputs
computed for the length $L_{\rm band}$ remain valid for every block of the band,
because every count inflation, tensor variance and row coefficient is nondecreasing in
the length. The root inputs are valid on $\Blk$ by
\cref{lem:record:monotone-endpoint} below $2^{52}$ (\cref{cert:upper-root-caps}) and by
\cref{lem:upper-completed-root} from $2^{52}$ on (\cref{cert:upper-completed-root}). The
weights satisfy \eqref{eq:upper-slice} and \eqref{eq:independent-fold-domination}
(\cref{cert:upper-completion}). Apply (i).
\end{proof}

Part (i) uses nothing about $\Blk$ except the validity of the inputs. It therefore
applies to a single prime $p\ge\Pfour$ of either class with the tail datum
(\cref{def:record:upper-bands}(e)), whose inputs are valid there (\cref{lem:tail-law}).
The data of the last band $[2^{64},2^{512})$ are class-three data and are not used
beyond $2^{512}$.

\begin{certificate}[Completion and domination gates]\label{cert:upper-completion}
At the baseline checkpoint, in each of the $310$ bands the $105$ inequalities \eqref{eq:upper-slice} (forty for
$w$, forty for $v$) and \eqref{eq:independent-fold-domination} (twenty-five) hold
exactly with the recorded inputs, $32550$ in all, and the weights are positive with
$w_{\rm LL}=1$. The record is
\path{research/further_2026_10_25/upper/upper_dense_refined.json}, with the fields
\texttt{inputs}, \texttt{reserved\_inputs}, \texttt{reserved\_state\_caps},
\texttt{weights}, \texttt{reserved\_weights} and \texttt{gates} for each band. Its
inputs, state caps and weights are those of two earlier records of that directory,
\path{upper_special.json} for the $301$ bands below $2^{52}$ and
\path{upper_endpoint_allowances.json} for the nine bands from $2^{52}$; the later
stages within that checkpoint change only the comparison ledgers
(\cref{sec:upper-comparison}). The final cross-weighted inputs and weights
are instead those of \cref{subsec:upper-cross-weighted}, with the same
completion and domination requirements. The
independent reviews of that directory re-check the gates on these records:
\path{special_upper_review.json} for the $301$ bands below $2^{52}$, whose root entries
and weights are fixed with the root caps of \cref{cert:upper-root-caps} ($24080$ slice
inequalities, $7525$ dominations and $903$ polynomials), and
\path{endpoint_allowances_review.json} for all $310$ bands of the allowance stage
($2790$ primitive prices, $24800$ slice inequalities, $7750$ dominations and $930$
physical polynomials), whose data the nine bands from $2^{52}$ keep. The script
\path{research/further_2026_10_25/upper/verify.py} replays the whole finite upper
extension from the certified baseline
\path{research/advance_2026_10_24/verification/selected.json}, without an optimizer,
and requires every regenerated record to equal its archived file; the selected
verification record \path{research/further_2026_10_25/verification/selected.json}
reports \texttt{PASS\_COMPLETE\_FURTHER\_UPPER\_EXTENSION}.

The $105$ gates of the tail datum are those of the band $[2^{64},2^{512})$ of
\path{research/roadmap_2026_10_18/item12/upper_v2_local.json}
(\texttt{PASS\_105\_EXACT\_COMPLETION\_AND\_DOMINATION\_GATES}), written by
\path{assemble_upper_v2.py} with the completion routine of
\path{assemble_upper_candidate.py} (same directory) and re-checked by
\path{research/push_2026_10_24/tools/check_upper.py} into
\path{research/push_2026_10_24/artifacts/coordinator/push_best_upper_review.json}
(\texttt{PASS\_INDEPENDENT\_GENERATED\_DIRECT\_UPPER\_PRICES\_AND\_GATES}: $1260$
gates, $105$ for each of the twelve bands of \path{upper_v2_local.json} from $2^{49}$
on, the band $[2^{64},2^{512})$ among them).
\end{certificate}

\paragraph{Consequences and what to read next.}
From Part~II: the slot array and its cuts (\cref{def:slot-array,lem:slot-array}),
square-affine invariance (\cref{lem:marginals}), admissible majorants
(\cref{lem:majorant}), admissibility (\cref{def:admissible}). From Part~IV: labellings
and counts (\cref{def:graphs:labeling,lem:extension,lem:four-count}), the two-root
identity and Weil's bound (\eqref{eq:two-root}, \cref{thm:weil}); spectral data and cut
kernels (\cref{def:spec:data,lem:spec:cutkernel}), block lifting
(\cref{lem:spec:tensorcomp}), the two-sided constant block
(\cref{prop:matching:positive}(a$'$)), pinned influence
(\cref{prop:middle-negative,cor:influence-schur}), recentering and restriction
(\cref{lem:recentering,lem:restriction}), row comparison (\cref{lem:row-comparison}); the
two spectral routes (\cref{thm:spectral-criterion,cor:spectral-criterion-blocks,thm:record:direct-star})
with the character expansion, the stars and the centring
(\cref{lem:spec:expansion,lem:spec:stars,prop:spec2:mucentre}), the frame inputs
(\cref{lem:record:frame-inputs}), the coupling block and its refinements
(\cref{cor:spec2:coupling,lem:push:localvariance,lem:projection-diagonal,lem:push:conditionalgram,lem:spec2:parity-transfer})
and the monotone spectral quantities (\cref{lem:spectral-monotone}); the complete actual
stars (\cref{lem:general-star,cor:star-matrix,lem:hidden-feature},
\cref{rem:realified-gates}); the two-variable sums and the moving-root moments
(\cref{thm:dl:H2,cor:moving-root-uniform,cor:moving-root-conditional}), the ordered star
marginal and the star mean (\cref{prop:advance:ordered-star,lem:star-mean}); the
compressed partition, the extraction bounds with the graded error pairs and the special
values, and the coincidence bounds
(\cref{lem:advance:compressed-partition,def:extraction:records,cor:extraction-leaf,cor:extraction-compressed,prop:pinned-extraction,lem:oct24-collision,prop:pinned-collision},
\eqref{eq:leaf-factors}); the log-convex nonstar coefficients
(\cref{lem:logconvex-coefficients}); the pinned conditional stars
(\cref{prop:conditional-star,rem:pinned-instances}, with \eqref{eq:row-anova-actual});
the joint rows (\cref{lem:joint-rows,cor:joint-row-split,lem:joint-potential}); and the
five-type completion with its joint form, its free allowance and its state caps
(\cref{lem:five-type}(b)--(f), \cref{cor:joint-completion-blocks,rem:five-type-fold-data}).

\paragraph{Hypotheses verified in this section: (CP1) for upper atoms.}
Every upper block carries the probability measure $\mu_\Blk$, reflection positive at
every cut of $\mathfrak F$ (the chamber cuts of both coordinates and $\iota$, not only the
seven sorting cuts), invariant under the square-affine maps, with uniform one-slot
marginals, and admissible (\cref{lem:upper-reference-rp}); with the band inputs of
\cref{def:record:upper-bands} its spectral caps and rows are valid on every block. The
tail datum (\cref{def:record:upper-bands}(e)) carries the same conclusions to single
primes $p\ge\Pfour$ of both classes (\cref{lem:tail-law}). This is the upper row of
(CP1) in \cref{tab:hypotheses}; the comparison law and its prices (CP4) are
\cref{prop:record:upper}.
